% =====================================================================
%  Relatório das Atividades Práticas de Fundamentos de DevOps
%  Projeto LaTeX para o Overleaf (compilador: pdfLaTeX)
% =====================================================================
\documentclass[12pt,a4paper,oneside]{article}

% =====================================================================
%  Preâmbulo: pacotes, página, cores, títulos e estilos de código
%  (compatível com pdfLaTeX no Overleaf)
% =====================================================================

% ---------- Codificação, idioma e fontes ----------
\usepackage[utf8]{inputenc}
\usepackage[T1]{fontenc}
\usepackage[brazil]{babel}
\usepackage{lmodern}
\usepackage{mathptmx}                 % Times no texto (padrão comum em trabalhos acadêmicos)
\renewcommand{\ttdefault}{lmtt}       % Latin Modern Mono para comandos e códigos
\usepackage{textcomp}
\usepackage{microtype}

% ---------- Página: margens no padrão ABNT ----------
\usepackage[a4paper,top=3cm,left=3cm,right=2cm,bottom=2cm,
            headheight=15pt,headsep=0.7cm,footskip=1cm]{geometry}
\usepackage{setspace}
\usepackage{indentfirst}
\setlength{\parindent}{1.25cm}
\setlength{\parskip}{0.2em}

% ---------- Figuras, tabelas, listas ----------
\usepackage{graphicx}
\usepackage{float}
\usepackage[table,dvipsnames]{xcolor}
\usepackage{booktabs}
\usepackage{tabularx}
\usepackage{array}
\usepackage{multirow}
\usepackage{longtable}
\usepackage{caption}
\captionsetup{font=small,labelfont=bf,labelsep=endash,justification=centering,
              singlelinecheck=false,skip=4pt}
\usepackage{enumitem}
\setlist{itemsep=0.1em,topsep=0.35em,parsep=0.1em}
\usepackage{dirtree}
\usepackage{tikz}
\usetikzlibrary{arrows.meta,positioning,fit,calc,shapes.geometric,backgrounds}
\usepackage{tcolorbox}
\tcbuselibrary{skins,breakable}
\usepackage{listings}
\usepackage{fancyhdr}
\usepackage{titlesec}
\usepackage{etoolbox}
\usepackage{needspace}
\usepackage{xurl}
\usepackage[round,authoryear]{natbib}     % citações autor-data: (AUTOR, ano)
\setcitestyle{semicolon,aysep={,},yysep={;}}
\usepackage[hidelinks,bookmarksnumbered,pdfusetitle]{hyperref}

% ---------- Cores do documento ----------
\definecolor{corprimaria}{HTML}{0B4A7D}   % azul
\definecolor{cordestaque}{HTML}{D9661F}   % laranja
\definecolor{corcinza}{HTML}{5B6670}
\definecolor{corverde}{HTML}{2E7D32}
\definecolor{corvermelha}{HTML}{B3261E}

% ---------- Títulos das seções ----------
\setcounter{secnumdepth}{3}
\setcounter{tocdepth}{2}
\titleformat{\section}{\normalfont\large\bfseries\color{corprimaria}}{\thesection}{0.6em}{\MakeUppercase}
\titleformat{\subsection}{\normalfont\normalsize\bfseries\color{corprimaria}}{\thesubsection}{0.6em}{}
\titleformat{\subsubsection}{\normalfont\normalsize\bfseries}{\thesubsubsection}{0.6em}{}
\titlespacing*{\section}{0pt}{1.2em}{0.8em}
\titlespacing*{\subsection}{0pt}{1.1em}{0.5em}
\titlespacing*{\subsubsection}{0pt}{0.9em}{0.4em}

% ---------- Cabeçalho e rodapé ----------
\pagestyle{fancy}
\fancyhf{}
\fancyhead[L]{\footnotesize\color{corcinza}Relatório das Atividades Práticas -- Fundamentos de DevOps}
\fancyhead[R]{\footnotesize\color{corcinza}\thepage}
\renewcommand{\headrulewidth}{0.4pt}
\renewcommand{\headrule}{\hbox to\headwidth{\color{gray!40}\leaders\hrule height \headrulewidth\hfill}}
\fancypagestyle{plain}{\fancyhf{}\fancyhead[R]{\footnotesize\color{corcinza}\thepage}\renewcommand{\headrulewidth}{0pt}}

% ---------- Estilos de código (listings) ----------
\definecolor{codfundo}{HTML}{F6F8FA}
\definecolor{codborda}{HTML}{D0D7DE}
\definecolor{codcomentario}{HTML}{6A737D}
\definecolor{codpalavra}{HTML}{0550AE}
\definecolor{codtexto}{HTML}{0A6E31}
\definecolor{codnumero}{HTML}{8C959F}
\definecolor{codsecao}{HTML}{8250DF}

\lstset{
  basicstyle=\ttfamily\footnotesize\setstretch{1.0},
  backgroundcolor=\color{codfundo},
  frame=single, framerule=0.4pt, rulecolor=\color{codborda},
  framesep=4pt, xleftmargin=2.4em, framexleftmargin=0pt,
  numbers=left, numberstyle=\tiny\color{codnumero}, numbersep=9pt,
  breaklines=true, breakatwhitespace=false,
  postbreak=\mbox{\textcolor{codnumero}{$\hookrightarrow$}\space},
  columns=fullflexible, keepspaces=true, showstringspaces=false,
  tabsize=2, upquote=true, extendedchars=true,
  commentstyle=\color{codcomentario}\itshape,
  keywordstyle=\color{codpalavra}\bfseries,
  stringstyle=\color{codtexto},
  captionpos=t, abovecaptionskip=2pt, belowcaptionskip=4pt,
  aboveskip=0.7em, belowskip=0.7em,
  literate=
    {á}{{\'a}}1 {é}{{\'e}}1 {í}{{\'i}}1 {ó}{{\'o}}1 {ú}{{\'u}}1
    {Á}{{\'A}}1 {É}{{\'E}}1 {Í}{{\'I}}1 {Ó}{{\'O}}1 {Ú}{{\'U}}1
    {à}{{\`a}}1 {À}{{\`A}}1
    {â}{{\^a}}1 {ê}{{\^e}}1 {ô}{{\^o}}1 {Â}{{\^A}}1 {Ê}{{\^E}}1 {Ô}{{\^O}}1
    {ã}{{\~a}}1 {õ}{{\~o}}1 {Ã}{{\~A}}1 {Õ}{{\~O}}1
    {ç}{{\c c}}1 {Ç}{{\c C}}1 {ü}{{\"u}}1
    {º}{{\textordmasculine}}1 {ª}{{\textordfeminine}}1
}

\lstdefinelanguage{yaml}{
  keywords={true,false,yes,no,null},
  sensitive=false,
  comment=[l]{\#},
  morestring=[b]',
  morestring=[b]",
}
\lstdefinelanguage{groovy}{
  morekeywords={pipeline,agent,any,stages,stage,steps,sh,echo,environment,options,
    triggers,post,success,failure,always,withCredentials,usernamePassword,string,
    git,writeFile,retry,sleep,pollSCM,buildDiscarder,logRotator,disableConcurrentBuilds},
  sensitive=true,
  morecomment=[l]{//},
  morecomment=[s]{/*}{*/},
  morestring=[b]",
  morestring=[b]',
}
\lstdefinelanguage{docker}{
  keywords={FROM,AS,RUN,CMD,LABEL,EXPOSE,ENV,ADD,COPY,ENTRYPOINT,VOLUME,USER,WORKDIR,ARG},
  sensitive=true,
  comment=[l]{\#},
  morestring=[b]",
}
\lstdefinelanguage{ini}{
  sensitive=false,
  morecomment=[l]{\#},
  morecomment=[l]{;},
  moredelim=[s][\color{codsecao}\bfseries]{[}{]},
}
\lstdefinestyle{terminal}{
  language=bash, numbers=none, xleftmargin=0pt,
  backgroundcolor=\color{codfundo},
}

\AtBeginDocument{%
  \renewcommand{\lstlistingname}{Código}%
  \renewcommand{\lstlistlistingname}{Lista de Códigos}%
}

% ---------- Estrutura de diretórios (dirtree) ----------
\renewcommand*\DTstyle{\ttfamily\small}
\setlength{\DTbaselineskip}{13pt}

% =====================================================================
%  Comandos e ambientes personalizados do relatório
% =====================================================================

% ---------- Fonte abaixo de figuras e tabelas (padrão ABNT) ----------
\newcommand{\fonte}[1]{\par\vspace{3pt}{\footnotesize Fonte: #1}\par}

% ---------- Texto a ser completado pelo aluno ----------
% Aparece destacado em amarelo no PDF para não passar despercebido.
\newcommand{\preencher}[1]{\colorbox{yellow!45}{\strut\footnotesize\bfseries[PREENCHER: #1]}}

% ---------- Código/comando curto dentro do texto ----------
\newcommand{\cod}{\lstinline[basicstyle=\ttfamily\small,keywordstyle=,stringstyle=,commentstyle=,
  postbreak={},breakindent=0pt]}

% ---------- Caixa exibida quando o print ainda não foi adicionado ----------
\newcommand{\caixaprint}[2]{%
  \begin{tcolorbox}[enhanced, width=0.86\linewidth, colback=gray!4,
      frame hidden, borderline={0.7pt}{0pt}{gray!55,dashed}, arc=1.5mm,
      halign=center, valign=center, height=3.6cm, fontupper=\small,
      before skip=2pt, after skip=2pt]
    {\color{cordestaque}\bfseries INSERIR PRINT AQUI}\par\smallskip
    {\ttfamily\footnotesize Salvar como: prints/#1.png}\par\medskip
    \begin{minipage}{0.92\linewidth}\centering\footnotesize #2\end{minipage}
  \end{tcolorbox}}

% ---------- Print de tela ----------
% \printtela[largura]{arquivo}{legenda}{o que capturar}
% Se existir prints/<arquivo>.png (ou .jpg/.jpeg), a imagem é inserida;
% caso contrário, aparece a caixa indicando qual print deve ser feito.
\newcommand{\printtela}[4][0.9]{%
  \begin{figure}[H]
    \centering
    \caption{#3}\label{fig:#2}
    \IfFileExists{prints/#2.png}{\includegraphics[width=#1\linewidth]{prints/#2.png}}{%
      \IfFileExists{prints/#2.jpg}{\includegraphics[width=#1\linewidth]{prints/#2.jpg}}{%
        \IfFileExists{prints/#2.jpeg}{\includegraphics[width=#1\linewidth]{prints/#2.jpeg}}{%
          \caixaprint{#2}{#4}}}}
    \fonte{Elaborado pelo autor (2026).}
  \end{figure}}

% ---------- Código-fonte a partir de arquivo do laboratório ----------
% \codigo[opções]{linguagem}{arquivo dentro de codigos/}{legenda}{rótulo}
%   -> listagem curta: fica inteira na mesma página (dentro de uma minipage)
% \codigolongo[opções]{linguagem}{arquivo}{legenda}{rótulo}
%   -> listagem longa: pode continuar na página seguinte
\newcommand{\codigo}[5][]{%
  \par\smallskip\noindent\begin{minipage}{\linewidth}%
  \lstinputlisting[language=#2,caption={[{#4}]{#4 -- \nolinkurl{#3}}},label={cod:#5},#1]{codigos/#3}%
  \end{minipage}\par\smallskip}
\newcommand{\codigolongo}[5][]{%
  \par\needspace{12\baselineskip}%
  \lstinputlisting[language=#2,caption={[{#4}]{#4 -- \nolinkurl{#3}}},label={cod:#5},#1]{codigos/#3}}

% ---------- Listagens digitadas no texto (não quebram entre páginas) ----------
% \begin{terminal}[caption={...},label={cod:...}] ... \end{terminal}  -> comandos de terminal
% \begin{listagem}[language=...,caption={...}] ... \end{listagem}    -> outros trechos de código
\lstnewenvironment{terminal}[1][]%
  {\endgraf\smallskip\noindent\minipage{\linewidth}\lstset{style=terminal,#1}}%
  {\endminipage\endgraf\smallskip}
\lstnewenvironment{listagem}[1][]%
  {\endgraf\smallskip\noindent\minipage{\linewidth}\lstset{#1}}%
  {\endminipage\endgraf\smallskip}

% ---------- Caixas de destaque ----------
\newtcolorbox{nota}[1][Observação]{enhanced, breakable, colback=corprimaria!4,
  colframe=corprimaria, boxrule=0pt, leftrule=2.5pt, arc=0pt, outer arc=0pt,
  left=7pt, right=7pt, top=5pt, bottom=5pt, fontupper=\small,
  before upper={\textbf{\color{corprimaria}#1:}\ }, before skip=8pt, after skip=8pt}

\newtcolorbox{atencao}[1][Atenção]{enhanced, breakable, colback=cordestaque!6,
  colframe=cordestaque, boxrule=0pt, leftrule=2.5pt, arc=0pt, outer arc=0pt,
  left=7pt, right=7pt, top=5pt, bottom=5pt, fontupper=\small,
  before upper={\textbf{\color{cordestaque}#1:}\ }, before skip=8pt, after skip=8pt}

% ---------- Estilos TikZ dos diagramas ----------
\tikzset{
  caixa/.style={draw=corprimaria, fill=corprimaria!6, rounded corners=3pt, align=center,
                font=\footnotesize, inner sep=5pt, line width=0.6pt},
  caixaescura/.style={draw=corprimaria, fill=corprimaria, text=white, rounded corners=3pt,
                align=center, font=\footnotesize\bfseries, inner sep=5pt},
  caixalaranja/.style={draw=cordestaque, fill=cordestaque!8, rounded corners=3pt, align=center,
                font=\footnotesize, inner sep=5pt, line width=0.6pt},
  caixacinza/.style={draw=gray!60, fill=gray!8, rounded corners=3pt, align=center,
                font=\footnotesize, inner sep=5pt},
  seta/.style={-{Stealth[length=2.4mm]}, line width=0.7pt, color=corcinza},
  setadupla/.style={{Stealth[length=2.4mm]}-{Stealth[length=2.4mm]}, line width=0.7pt, color=corcinza},
  rotulo/.style={font=\scriptsize, text=corcinza, fill=white, inner sep=1.5pt},
}


% ---------- Dados do trabalho (confira e ajuste) ----------
\newcommand{\autor}{Mateus Zanatta Mariani}
\newcommand{\instituicao}{Centro Universitário Senac RS -- UNISENAC Pelotas}
\newcommand{\curso}{Curso Superior de Tecnologia em Análise e Desenvolvimento de Sistemas}
\newcommand{\disciplina}{Fundamentos de DevOps}
\newcommand{\professor}{Prof. Augusto da Rosa Muniz}
\newcommand{\titulo}{Relatório das Atividades Práticas de Fundamentos de DevOps}
\newcommand{\subtitulo}{Vagrant, Git e GitHub, Ansible, Docker, SonarQube e Jenkins}
\newcommand{\local}{Pelotas}
\newcommand{\ano}{2026}
\newcommand{\repositorio}{https://github.com/mateus-zanatta/trabalho-devops-2}

\hypersetup{pdftitle={\titulo}, pdfauthor={\autor}, pdfsubject={\disciplina}}

\begin{document}

% ---------------------------------------------------------------------
%  Elementos pré-textuais
% ---------------------------------------------------------------------
% =====================================================================
%  Capa, folha de rosto, sumário e listas
% =====================================================================

% ---------- Capa ----------
\begin{titlepage}
  \centering
  {\large\bfseries \MakeUppercase{\instituicao}\par}
  \vspace{0.25cm}
  {\normalsize \MakeUppercase{\curso}\par}
  \vspace{0.25cm}
  {\normalsize Unidade Curricular: \disciplina\par}

  \vspace{3.2cm}
  {\large\bfseries \MakeUppercase{\autor}\par}

  \vspace{4cm}
  {\color{corprimaria}\rule{0.6\linewidth}{0.8pt}\par}
  \vspace{0.5cm}
  {\LARGE\bfseries\color{corprimaria} \titulo\par}
  \vspace{0.45cm}
  {\large \subtitulo\par}
  \vspace{0.5cm}
  {\color{corprimaria}\rule{0.6\linewidth}{0.8pt}\par}

  \vfill
  {\normalsize \local\par}
  {\normalsize \ano\par}
\end{titlepage}

% ---------- Folha de rosto ----------
\thispagestyle{empty}
\begin{center}
  {\large\bfseries \MakeUppercase{\autor}\par}

  \vspace{5cm}
  {\Large\bfseries \titulo\par}
  \vspace{0.4cm}
  {\large \subtitulo\par}
\end{center}

\vspace{2cm}
\noindent\hfill
\begin{minipage}{0.55\linewidth}
  \small\setstretch{1.0}
  Relatório técnico apresentado à Unidade Curricular de \disciplina{}, do \curso{} do
  \instituicao{}, como requisito parcial de avaliação.

  \vspace{0.6cm}
  Professor: \professor{}.
\end{minipage}

\vfill
\begin{center}
  {\normalsize \local\par}
  {\normalsize \ano\par}
\end{center}
\clearpage

% ---------- Sumário e listas ----------
\pagestyle{empty}
{\hypersetup{linkcolor=black}
 \setstretch{1.15}
 \tableofcontents
 \thispagestyle{empty}
 \clearpage
 \listoffigures
 \thispagestyle{empty}
 \clearpage
 \lstlistoflistings
 \thispagestyle{empty}
 \clearpage
}


\onehalfspacing

% ---------------------------------------------------------------------
%  Elementos textuais
% ---------------------------------------------------------------------
\pagestyle{fancy}
% =====================================================================
%  1 Introdução
% =====================================================================
\section{Introdução}

Durante muito tempo, a entrega de software foi organizada em equipes separadas: os desenvolvedores
escreviam o código e as equipes de operações (infraestrutura) eram responsáveis por instalá-lo e
mantê-lo funcionando em servidores físicos, sujeitos a falhas. Com o modelo em cascata
(\textit{waterfall}), as entregas eram lentas, cada \textit{deploy} se tornava um ``super evento''
que envolvia vários times isolados em silos, e boa parte da configuração de servidores ainda era
feita de forma manual \citep{muniz2026}. O resultado era pouca inovação, retrabalho e a famosa
frase ``funciona na minha máquina''.

O \textbf{DevOps} surgiu como resposta a esse cenário. Ele pode ser definido como um conjunto de
práticas que combina o desenvolvimento de software (\textit{Dev}) e as operações de TI (\textit{Ops})
com o objetivo de reduzir o ciclo de desenvolvimento, oferecendo entrega contínua de software de
alta qualidade \citep{muniz2026,kim2018}. É importante destacar que DevOps não é uma ferramenta,
e sim uma cultura e uma metodologia de trabalho: desenvolvedores e operações passam a trabalhar
juntos durante todo o ciclo de vida do software, e o time de operações adota técnicas típicas do
desenvolvimento, como versionamento, testes e automação. A Figura~\ref{fig:ciclo-devops} resume as
etapas desse ciclo, que se repete continuamente.

\begin{figure}[H]
  \centering
  \caption{Etapas do ciclo DevOps}\label{fig:ciclo-devops}
  \begin{tikzpicture}
    \foreach \nome/\cor [count=\i from 0] in {Definição/corprimaria, Planejar/corprimaria,
        Codificar/corprimaria, Build/corprimaria, Testar/corprimaria, Liberar/cordestaque,
        Deploy/cordestaque, Operar/cordestaque, Monitorar/cordestaque} {
      \node[draw=\cor, fill=\cor!10, rounded corners=3pt, minimum width=2.15cm,
            minimum height=0.72cm, font=\footnotesize\bfseries, text=\cor!85!black]
            (n\i) at ({4.7*cos(90-\i*40)}, {2.85*sin(90-\i*40)}) {\nome};
    }
    \foreach \i [evaluate=\i as \j using {int(mod(\i+1,9))}] in {0,...,8} {
      \draw[seta] (n\i) to[bend left=10] (n\j);
    }
    \node[font=\normalsize\bfseries, text=corprimaria] at (0,0.3) {DevOps};
    \node[align=center, font=\scriptsize, text=corcinza] at (0,-0.35)
         {ciclo contínuo de\\entrega e \textit{feedback}};
    \node[anchor=north, font=\scriptsize, text=corcinza] at (0,-3.55)
         {\textcolor{corprimaria}{\rule{1.4ex}{1.4ex}} etapas mais ligadas ao desenvolvimento (\textit{Dev})
          \qquad \textcolor{cordestaque}{\rule{1.4ex}{1.4ex}} etapas mais ligadas às operações (\textit{Ops})};
  \end{tikzpicture}
  \fonte{Adaptado de Muniz (2026), Aula 1.}
\end{figure}

A integração entre desenvolvimento e operações é importante porque elimina a ``passagem de bastão''
entre equipes. Quando as duas áreas compartilham responsabilidades, ferramentas e objetivos, os
problemas são descobertos mais cedo, as entregas ficam menores e mais frequentes, e o
\textit{feedback} sobre cada mudança chega rapidamente a quem a fez. Entre os benefícios
apresentados na disciplina estão o aumento da velocidade de entrega, a criação de uma cultura de
agilidade e automação, a aceleração dos \textit{releases}, o apoio à área de negócio
(\textit{time to market}), o aumento da segurança e da estabilidade e uma colaboração muito maior
entre as pessoas \citep{muniz2026}. Ainda assim, DevOps não é ``bala de prata'': ambientes pequenos,
muito controlados ou com poucas mudanças podem funcionar bem com modelos mais tradicionais.

Para que essa cultura funcione na prática, algumas práticas técnicas são fundamentais:

\begin{itemize}
  \item \textbf{Automação}: tarefas repetitivas (instalar pacotes, configurar serviços, compilar,
        testar, publicar) passam a ser executadas por scripts e ferramentas, o que reduz o erro
        humano e garante que o processo seja sempre igual.
  \item \textbf{Versionamento}: todo artefato do projeto -- código, scripts e configurações -- fica
        em um sistema de controle de versão como o Git, com histórico, rastreabilidade e a
        possibilidade de voltar a uma versão anterior.
  \item \textbf{Infraestrutura como Código (IaC)}: servidores e ambientes são descritos em arquivos
        de texto (como \texttt{Vagrantfile} e \textit{playbooks} do Ansible), que podem ser
        versionados, revisados e executados quantas vezes for necessário, sempre gerando o mesmo
        resultado.
  \item \textbf{Containers}: a aplicação é empacotada junto com todas as suas dependências,
        tornando-se portátil e isolada, o que elimina conflitos de versões entre ambientes.
  \item \textbf{Análise de código}: ferramentas de análise estática, como o SonarQube, apontam
        \textit{bugs}, vulnerabilidades e problemas de manutenção antes que o código chegue à produção.
  \item \textbf{Pipelines}: ferramentas como o Jenkins encadeiam automaticamente as etapas de
        \textit{build}, teste, análise e \textit{deploy} a cada alteração, implementando a
        integração e a entrega contínuas (CI/CD).
\end{itemize}

O objetivo das práticas realizadas na disciplina foi construir, passo a passo, um laboratório que
simula um ambiente corporativo real: máquinas virtuais criadas com o Vagrant, código e configurações
versionados com Git e GitHub, servidores de aplicação e de banco de dados configurados pelo Ansible,
a aplicação Java \textit{Easy Notes} empacotada em container com Docker, a qualidade do código
verificada pelo SonarQube e todo o fluxo automatizado em um \textit{pipeline} no Jenkins. Mais do que
aprender cada ferramenta isoladamente, a proposta foi entender como elas se conectam dentro do
ciclo DevOps.

Este relatório está organizado da seguinte forma: a Seção~\ref{sec:desenvolvimento} descreve o
ambiente do laboratório e cada atividade prática (Vagrant, Git e GitHub, Ansible, Docker, SonarQube
e Jenkins), além da integração entre as ferramentas; a Seção~\ref{sec:conclusao} apresenta a
conclusão; e a Seção~\ref{sec:reflexao} traz a resposta à pergunta reflexiva final.

% =====================================================================
%  2 Desenvolvimento -- visão geral do ambiente
% =====================================================================
\section{Desenvolvimento}\label{sec:desenvolvimento}

As práticas foram feitas em um computador com Windows, usando o \textbf{Vagrant} para criar as
máquinas virtuais e o \textbf{VMware Workstation} para executá-las. Todas as VMs usam o Rocky Linux 9
(\textit{box} \texttt{generic/rocky9}) e estão ligadas na rede privada \texttt{192.168.1.0/24}. A
Tabela~\ref{tab:maquinas} resume as máquinas criadas ao longo da disciplina. Os arquivos completos
do laboratório estão versionados no repositório do trabalho, na pasta \texttt{relatorio/codigos}
(\url{\repositorio}); senhas e chaves foram substituídas por marcadores.

\begin{table}[H]
  \centering
  \caption{Máquinas virtuais criadas no laboratório}\label{tab:maquinas}
  \footnotesize
  \setlength{\tabcolsep}{5pt}
  \renewcommand{\arraystretch}{1.3}
  \begin{tabularx}{\linewidth}{>{\ttfamily}l >{\ttfamily}l l c >{\raggedright\arraybackslash}X}
    \toprule
    \normalfont\bfseries Hostname & \normalfont\bfseries IP privado &
    \bfseries Porta redirecionada (VM $\rightarrow$ host) & \bfseries RAM/vCPU & \bfseries Função \\
    \midrule
    \normalfont(padrão) & \normalfont-- (só NAT) & 80 $\rightarrow$ 127.0.0.1:8080 & 1 GB / 2 & Laboratório 1: Apache servindo uma página HTML \\
    control-node & 192.168.1.2 & -- & 1 GB / 2 & \textit{Control Node} do Ansible \\
    app01 & 192.168.1.3 & 80 $\rightarrow$ 8080 & 1 GB / 2 & Aplicação Java \textit{Easy Notes} (porta 8080) \\
    db01 & 192.168.1.4 & 3306 $\rightarrow$ 3306 & 1 GB / 2 & Banco de dados MariaDB \\
    docker-notes & 192.168.1.5 & 8080 $\rightarrow$ 127.0.0.1:8082 & 2 GB / 2 & Docker: aplicação em container \\
    sonarqube & 192.168.1.6 & 9000 $\rightarrow$ 127.0.0.1:9000 & 4 GB / 2 & SonarQube e SonarScanner \\
    jenkins & 192.168.1.7 & 8080 $\rightarrow$ 127.0.0.1:8081 & 4 GB / 2 & Jenkins (CI/CD) \\
    \bottomrule
  \end{tabularx}
  \fonte{Elaborado pelo autor (2026).}
\end{table}

% =====================================================================
%  2.1 Atividade com Vagrant
% =====================================================================
\subsection{Atividade com Vagrant}\label{sec:vagrant}

\subsubsection{O que é o Vagrant e para que foi usado}

O \textbf{Vagrant}, da HashiCorp, é uma ferramenta de linha de comando que cria e gerencia máquinas
virtuais a partir de um arquivo de texto, o \texttt{Vagrantfile}. Nele são definidos a
\textbf{\textit{box}} (imagem base do sistema operacional), a rede, as portas, as pastas compartilhadas
e os scripts de \textbf{provisionamento}, executados automaticamente na criação da VM; quem executa a
máquina é o \textbf{\textit{provider}} (VirtualBox, VMware, Hyper-V\ldots) \citep{hashicorp2026}.
Com isso, o ambiente deixa de ser montado ``na mão'' e passa a ser descrito por código, o que acaba
com o problema do ``funciona na minha máquina''.

Na disciplina, o Vagrant foi a base de todo o laboratório. Na primeira prática, ele criou uma VM com o
Apache instalado automaticamente, servindo uma página HTML acessada pelo navegador do Windows em
\url{http://localhost:8080}. Depois, o mesmo padrão -- um diretório com \texttt{Vagrantfile} e
\texttt{provision.sh} para cada servidor -- foi usado para criar todas as VMs das práticas seguintes:

\begin{itemize}
  \item \textbf{Box}: \texttt{generic/rocky9} (Rocky Linux 9, compatível com o Red Hat).
  \item \textbf{Provider}: VirtualBox na primeira aula e, a partir da Aula 3, \textbf{VMware
        Workstation}, com o \textit{plugin} \path{vagrant-vmware-desktop}
        (\cod|vagrant plugin install vagrant-vmware-desktop|) e o serviço Vagrant VMware Utility.
  \item \textbf{Máquinas}: sete VMs ao longo da disciplina, com os \textit{hostnames}, IPs e portas
        redirecionadas da Tabela~\ref{tab:maquinas}. As portas do host ficaram acima de 1024 e sem
        repetição, para que várias VMs pudessem ficar ligadas ao mesmo tempo.
\end{itemize}

\subsubsection{Explicação do Vagrantfile}

O Código~\ref{cod:vagrantfile-app01} mostra o \texttt{Vagrantfile} do servidor \texttt{app01}; os das
outras VMs seguem o mesmo modelo, mudando nome, IP e portas.

\codigo{Ruby}{02-ansible-lab/app01/Vagrantfile}{Vagrantfile do servidor app01}{vagrantfile-app01}

\begin{itemize}
  \item \cod|Vagrant.configure("2")|: abre o bloco de configuração (versão 2 do formato).
  \item \cod|config.vm.box|: define a \textit{box}, baixada uma única vez e clonada para cada VM.
  \item \cod|config.vm.hostname|: define o nome da máquina.
  \item \cod|"private_network", ip: ...|: cria a rede privada com IP fixo, usada na comunicação
        entre as VMs.
  \item \cod|"forwarded_port", guest: 80, host: 8080|: redireciona a porta 8080 do Windows para a
        porta 80 da VM.
  \item \cod|config.vm.provision "shell", path: "provision.sh"|: executa o script de
        provisionamento no primeiro \cod|vagrant up|. No primeiro laboratório ele instalava o Apache
        e liberava a porta 80 no \textit{firewall}; no \texttt{app01} e no \texttt{db01}, ele
        autoriza a chave SSH do \textit{control-node} (Seção~\ref{sec:ansible}).
  \item \cod|config.vm.provider "vmware_desktop"|: ajusta, no VMware, a quantidade de memória
        (\texttt{memsize}) e de CPUs (\texttt{numvcpus}) da VM.
\end{itemize}

Além disso, \cod|config.vm.synced_folder ".", "/vagrant"| compartilha a pasta do projeto com a VM;
foi assim que a página HTML do primeiro laboratório e os arquivos do Ansible chegaram às máquinas.

\subsubsection{Comandos utilizados}

\begin{table}[H]
  \centering
  \caption{Comandos do Vagrant utilizados}\label{tab:comandos-vagrant}
  \small
  \renewcommand{\arraystretch}{1.15}
  \begin{tabularx}{\linewidth}{>{\ttfamily}l >{\raggedright\arraybackslash}X}
    \toprule
    \normalfont\bfseries Comando & \textbf{O que faz} \\
    \midrule
    vagrant init & Cria um \texttt{Vagrantfile} inicial no diretório atual. \\
    vagrant up & Cria e inicia a VM; na primeira vez, baixa a \textit{box} e executa o provisionamento. \\
    vagrant status & Mostra o estado da VM (\textit{running}, \textit{poweroff}, \textit{not created}). \\
    vagrant ssh & Abre um terminal dentro da VM. \\
    vagrant halt & Desliga a VM sem apagá-la. \\
    vagrant reload & Reinicia a VM aplicando mudanças do \texttt{Vagrantfile} (com \texttt{-{}-provision}, reexecuta o script). \\
    vagrant destroy -f & Apaga a VM (a \textit{box} continua salva) para recriar o ambiente do zero. \\
    \bottomrule
  \end{tabularx}
  \fonte{Elaborado pelo autor (2026), com base em Muniz (2026).}
\end{table}

Para validar o primeiro laboratório, dentro da VM foram usados \cod|curl localhost| e
\cod|sudo systemctl status httpd|, e no Windows o acesso a \url{http://localhost:8080}.

\printtela{vagrant-up}{Criação e provisionamento de uma VM com \texttt{vagrant up}}%
  {Terminal (PowerShell) com a saída do \texttt{vagrant up}: criação da VM com o provider
   \texttt{vmware\_desktop}, redirecionamento de porta e mensagens do \texttt{provision.sh}.}

\printtela{vagrant-status}{VMs do laboratório em execução}%
  {Terminal com \texttt{vagrant global-status} (ou \texttt{vagrant status}) mostrando as VMs no
   estado \textit{running}.}

\subsubsection{O papel do Vagrant na criação de ambientes reprodutíveis}

Como toda a configuração fica no \texttt{Vagrantfile} e no \texttt{provision.sh}, foi possível
destruir e recriar as máquinas várias vezes durante o semestre -- por exemplo, para reaplicar a chave
SSH no laboratório de Ansible -- e o resultado era sempre o mesmo. Isso torna o ambiente
\textbf{reprodutível}: qualquer pessoa que clone o repositório e execute \cod|vagrant up| obtém as
mesmas máquinas, com o mesmo sistema, IPs, portas e pacotes. Em um laboratório de DevOps, isso permite
errar, testar e voltar ao estado inicial sem medo, tratando a infraestrutura como código.

% =====================================================================
%  2.2 Atividade com Git e GitHub
% =====================================================================
\subsection{Atividade com Git e GitHub}\label{sec:git}

\subsubsection{O que é Git}

O \textbf{Git} é um sistema de controle de versão \textbf{distribuído}: ele gerencia as mudanças
feitas nos arquivos de um projeto e mantém um histórico completo de modificações, registrando
\textbf{quem} alterou, \textbf{quando} alterou e \textbf{o que} foi alterado \citep{chacon2014}.
Uma boa analogia, usada em aula, é a do \textit{snapshot} de máquina virtual: cada versão registrada
é uma ``fotografia'' do estado do projeto naquele momento, o que permite voltar no tempo caso um
erro seja introduzido \citep{muniz2026}. Com isso o Git oferece rastreabilidade (é possível saber
qual linha mudou e por quê), evita a perda de código e permite que várias pessoas trabalhem no mesmo
projeto sem sobrescrever o trabalho umas das outras. O \textbf{GitHub}, por sua vez, é uma plataforma
\textit{on-line} que hospeda repositórios Git remotos e acrescenta recursos de colaboração, como
\textit{pull requests}, \textit{issues} e automações.

\subsubsection{Repositório local e repositório remoto}

O \textbf{repositório local} fica no próprio computador, dentro da pasta oculta \texttt{.git}
criada pelo \cod|git init| (ou pelo \cod|git clone|). Ele guarda todo o histórico do projeto, de
modo que é possível criar versões, consultar o histórico e criar \textit{branches} mesmo sem internet.
Como foi lembrado em aula, ``o Git mora dentro da pasta \texttt{.git}: se ela some, o Git também''.
O \textbf{repositório remoto} é uma cópia hospedada em um servidor -- no nosso caso, o GitHub --,
identificada normalmente pelo nome \texttt{origin}. Ele funciona como ponto central de sincronização
da equipe e como \textit{backup}: como cada desenvolvedor tem uma cópia completa do repositório,
mesmo que o servidor fique indisponível o código não se perde. A Figura~\ref{fig:fluxo-git} mostra
as áreas do Git e os comandos que movem as alterações entre elas.

\begin{figure}[H]
  \centering
  \caption{Áreas do Git e comandos que movem as alterações entre elas}\label{fig:fluxo-git}
  \begin{tikzpicture}
    \tikzset{area/.style={caixa, text width=2.25cm, minimum height=1.25cm}}
    \node[area] (wd) {\textbf{Diretório de trabalho}\\{\scriptsize arquivos editados}};
    \node[area, right=1.65cm of wd] (st) {\textbf{Área de preparação}\\{\scriptsize \textit{staging} (índice)}};
    \node[area, right=1.65cm of st] (lr) {\textbf{Repositório local}\\{\scriptsize pasta \texttt{.git}}};
    \node[caixalaranja, text width=2.25cm, minimum height=1.25cm, right=1.65cm of lr] (rr)
         {\textbf{Repositório remoto}\\{\scriptsize GitHub (\texttt{origin})}};
    \draw[seta] (wd.east) -- node[rotulo, above=1pt] {\texttt{git add}} (st.west);
    \draw[seta] (st.east) -- node[rotulo, above=1pt] {\texttt{git commit}} (lr.west);
    \draw[seta] ([yshift=4pt]lr.east) -- node[rotulo, above=1pt] {\texttt{git push}} ([yshift=4pt]rr.west);
    \draw[seta] ([yshift=-4pt]rr.west) -- node[rotulo, below=1pt] {\texttt{git fetch}} ([yshift=-4pt]lr.east);
    \draw[seta] (rr.south) -- ++(0,-0.75) -| node[rotulo, pos=0.25, below=1pt]
         {\texttt{git pull} / \texttt{git clone}: trazem as alterações do remoto para o diretório de trabalho} (wd.south);
    \begin{scope}[on background layer]
      \node[draw=corprimaria!45, dashed, rounded corners=4pt, fit=(wd)(st)(lr), inner sep=7pt,
            label={[font=\scriptsize\bfseries, text=corprimaria]above:Computador do desenvolvedor (repositório local)}] {};
    \end{scope}
  \end{tikzpicture}
  \fonte{Elaborado pelo autor (2026), com base em Chacon e Straub (2014).}
\end{figure}

\subsubsection{Commit, push, pull e branch}

\begin{itemize}
  \item \textbf{Commit}: registra no repositório local uma nova versão com as alterações que estavam
        na área de preparação. Cada \textit{commit} tem um identificador único (\textit{hash}), autor,
        data e uma mensagem que explica o que foi feito, por exemplo
        \cod|git commit -m "Adiciona playbook do app01"|.
  \item \textbf{Push}: envia os \textit{commits} do repositório local para o repositório remoto
        (\cod|git push|), tornando-os visíveis para a equipe.
  \item \textbf{Pull}: busca as alterações do repositório remoto e as integra à \textit{branch} local
        (\cod|git pull| equivale a um \cod|git fetch| seguido de um \cod|git merge|).
  \item \textbf{Branch}: é uma ramificação, ou seja, uma linha de desenvolvimento separada da principal
        (\texttt{main}). Permite desenvolver uma funcionalidade ou corrigir um erro de forma isolada,
        com \textit{commits} próprios, sem quebrar o código principal. Quando o trabalho termina, a
        \textit{branch} é integrada de volta com \cod|git merge|.
\end{itemize}

Na disciplina também foi apresentado o \textbf{Git Flow}, um modelo de organização de
\textit{branches} com cinco tipos: \texttt{main} (versões estáveis em produção, marcadas com
\textit{tags} como \texttt{v1.0}), \texttt{develop} (integração do desenvolvimento),
\texttt{feature} (novas funcionalidades, criadas a partir da \texttt{develop}), \texttt{release}
(preparação de uma nova versão, com testes finais) e \texttt{hotfix} (correções urgentes que nascem
da \texttt{main} e voltam para a \texttt{main} e para a \texttt{develop}). O Git Flow é indicado para
projetos grandes, com vários ambientes e lançamentos programados; já o \textit{feature branching},
mais simples, combina melhor com projetos menores e com integração contínua \citep{muniz2026}.

\subsubsection{Como o Git foi utilizado nas atividades}

Primeiro, o Git foi instalado e configurado com a identidade do autor, que fica gravada em cada
\textit{commit}. Em seguida, foi criado um repositório no GitHub e o projeto local foi inicializado
e ligado a ele, seguindo a sequência do Código~\ref{cod:git-sequencia}. Antes do primeiro
\textit{commit} foi criado o arquivo \texttt{.gitignore} (Código~\ref{cod:gitignore}), para que
arquivos que não devem ser publicados -- como a pasta \texttt{.vagrant/}, que guarda o estado
local das VMs, logs e arquivos \texttt{.env} com segredos -- ficassem fora do repositório.

\begin{terminal}[caption={Sequência de comandos usada para versionar o laboratório no GitHub},label={cod:git-sequencia}]
git --version                                   # confirma a instalação
git config --global user.name  "Seu Nome"       # identidade gravada nos commits
git config --global user.email "seu@email.com"
git config --global --list                      # confere a configuração

cd lab-devops                                   # ATENÇÃO: dentro da pasta do projeto
git init                                        # cria o repositório local (.git)
echo "# Lab DevOps" > README.md
git status                                      # mostra arquivos novos/alterados
git add .                                       # adiciona tudo à área de preparação
git commit -m "Primeiro commit do projeto"      # registra a primeira versão
git branch -M main                              # renomeia a branch principal para main
git remote add origin https://github.com/usuario/lab-devops.git
git remote -v                                   # confere o repositório remoto
git push -u origin main                         # envia os commits para o GitHub
\end{terminal}

% Conteúdo do arquivo codigos/.gitignore (inserido aqui porque o Overleaf pode ignorar arquivos ocultos)
\begin{listagem}[language=ini,caption={[Arquivo .gitignore do laboratório]{Arquivo .gitignore do laboratório -- \nolinkurl{.gitignore}}},label={cod:gitignore}]
# Arquivos que NÃO devem ir para o GitHub
.vagrant/
*.log
.env
.DS_Store
\end{listagem}

A partir daí, a cada etapa concluída (por exemplo, um novo \textit{playbook} ou um ajuste no
\texttt{Vagrantfile}) o ciclo \cod|git status| $\rightarrow$ \cod|git add| $\rightarrow$
\cod|git commit| $\rightarrow$ \cod|git push| foi repetido. Para alterações maiores foi usada uma
\textit{branch} separada, integrada à \texttt{main} somente depois de testada
(Código~\ref{cod:git-branch}). Este próprio relatório, junto com todos os arquivos do laboratório,
está versionado no repositório \url{\repositorio}.

\begin{terminal}[caption={Uso de branch para desenvolver uma alteração isolada},label={cod:git-branch}]
git checkout -b feature/playbook-app            # cria a branch e muda para ela
git add ansible/playbooks/app.yml
git commit -m "Adiciona playbook de configuração do app01"
git checkout main                               # volta para a branch principal
git pull                                        # garante que a main está atualizada
git merge feature/playbook-app                  # integra a alteração
git push                                        # publica no GitHub
git log --oneline --graph                       # histórico resumido
\end{terminal}

\subsubsection{Comandos utilizados}

\begin{table}[H]
  \centering
  \caption{Comandos do Git utilizados}\label{tab:comandos-git}
  \small
  \renewcommand{\arraystretch}{1.2}
  \begin{tabularx}{\linewidth}{>{\ttfamily}l >{\raggedright\arraybackslash}X}
    \toprule
    \normalfont\bfseries Comando & \textbf{Função} \\
    \midrule
    git init & Cria um repositório Git na pasta atual. \\
    git clone URL & Copia um repositório remoto (com todo o histórico) para o computador. \\
    git status & Mostra arquivos novos, alterados e preparados para \textit{commit}. \\
    git add . & Adiciona as alterações à área de preparação (\textit{staging}). \\
    git commit -m "msg" & Salva uma nova versão no repositório local. \\
    git remote add origin URL & Liga o repositório local ao remoto (GitHub). \\
    git pull & Busca e integra as alterações do remoto. \\
    git push & Envia os \textit{commits} locais para o remoto. \\
    git branch & Lista ou cria \textit{branches}. \\
    git checkout -b nome & Cria uma \textit{branch} e muda para ela. \\
    git merge nome & Integra a \textit{branch} indicada na \textit{branch} atual. \\
    git log -{}-oneline & Mostra o histórico de \textit{commits} de forma resumida. \\
    \bottomrule
  \end{tabularx}
  \fonte{Elaborado pelo autor (2026), com base em Muniz (2026).}
\end{table}

\subsubsection{Evidências da execução}

\printtela{github-repositorio}{Repositório do trabalho no GitHub}%
  {Página do repositório no GitHub mostrando os arquivos e pastas versionados
   (por exemplo, \texttt{relatorio/} e \texttt{README.md}) e o último \textit{commit}.}

\printtela{git-status}{Saída do comando \texttt{git status}}%
  {Terminal na pasta do projeto mostrando o \texttt{git status} (arquivos modificados/não rastreados
   ou ``nothing to commit, working tree clean'').}

\printtela{git-log}{Histórico de commits (\texttt{git log -{}-oneline})}%
  {Terminal com \texttt{git log -{}-oneline -{}-graph} mostrando os \textit{commits} realizados.}

\printtela{git-push}{Envio das alterações para o GitHub com \texttt{git push}}%
  {Terminal com a saída do \texttt{git push} (ou a aba \textit{Commits} do GitHub com o histórico).}

\printtela{github-arquivo}{Arquivo versionado visualizado no GitHub}%
  {GitHub exibindo um arquivo versionado, por exemplo o \texttt{Vagrantfile} ou o
   \texttt{playbooks/db.yml}, com o histórico de alterações.}

\subsubsection{Por que o versionamento é importante em DevOps}

No DevOps, o repositório Git é a ``fonte única da verdade'': tudo o que define o sistema -- o código
da aplicação, os \texttt{Vagrantfile}, os \textit{playbooks} do Ansible, o \texttt{Dockerfile} e o
\texttt{Jenkinsfile} -- fica versionado. Isso traz rastreabilidade (cada mudança tem autor, data e
motivo), segurança para experimentar (é possível voltar a uma versão que funcionava), colaboração
organizada por meio de \textit{branches} e revisões, e é a base da automação: os \textit{pipelines}
de CI/CD são disparados justamente por um \textit{commit} ou \textit{push} no repositório. Sem
versionamento, não existe Infraestrutura como Código nem integração contínua.

% =====================================================================
%  2.3 Atividade com Ansible
% =====================================================================
\subsection{Atividade com Ansible}\label{sec:ansible}

\subsubsection{O que é o Ansible}

O \textbf{Ansible} é uma ferramenta \textit{open source}, mantida pela Red Hat, para automação de
TI: gerenciamento de configuração, provisionamento de servidores, \textit{deploy} de aplicações e
orquestração de ambientes \citep{redhat2026}. Em vez de configurar máquina por máquina, o
administrador descreve em arquivos YAML (os \textit{playbooks}) o estado desejado dos servidores --
quais pacotes devem estar instalados, quais serviços devem estar rodando, quais arquivos devem
existir -- e o Ansible executa tudo automaticamente. Suas principais características são a
arquitetura \textbf{\textit{agentless}} (não é preciso instalar nenhum agente nos servidores
gerenciados), a comunicação via \textbf{SSH} no Linux (ou WinRM no Windows), a sintaxe simples em
YAML e a \textbf{idempotência}. O exemplo dado em aula resume bem a ideia: instalar o Java em dez
servidores manualmente é demorado e sujeito a erro; com o Ansible, basta um \textit{playbook} e um
único comando \citep{muniz2026}.

\subsubsection{Infraestrutura como Código}

\textbf{Infraestrutura como Código} (\textit{Infrastructure as Code}, IaC) é a prática de definir e
gerenciar servidores, redes e configurações por meio de arquivos de texto, em vez de processos
manuais. Esses arquivos podem ser versionados no Git, revisados, reutilizados e executados quantas
vezes for necessário, sempre produzindo o mesmo ambiente. No laboratório, a IaC apareceu em duas
camadas complementares: o \textbf{Vagrant} descreve e cria as máquinas virtuais
(\texttt{Vagrantfile}), e o \textbf{Ansible} descreve e aplica a configuração dentro delas
(\textit{playbooks} e \textit{roles}). Como destacado em aula: ``não estamos criando apenas pastas;
estamos simulando servidores reais dentro de um ambiente automatizado'' \citep{muniz2026}.

\subsubsection{Arquitetura: Control Node, Managed Hosts e inventário}

A arquitetura do Ansible é composta por poucos elementos \citep{muniz2026,redhat2026}:

\begin{itemize}
  \item \textbf{Control Node}: máquina onde o Ansible está instalado e de onde os comandos e
        \textit{playbooks} são executados. Ela armazena o inventário, os \textit{playbooks} e as
        chaves SSH. No laboratório, é o \texttt{control-node} (\texttt{192.168.1.2}), o ``cérebro''
        do ambiente.
  \item \textbf{Managed Hosts}: servidores gerenciados pelo Ansible. Não precisam de agente, apenas
        de SSH ativo e Python instalado. No laboratório, são o \texttt{app01} (\texttt{192.168.1.3},
        servidor de aplicação) e o \texttt{db01} (\texttt{192.168.1.4}, servidor de banco de dados).
  \item \textbf{Inventory} (inventário): arquivo que lista os servidores e os organiza em grupos.
  \item \textbf{Playbook}: arquivo YAML com uma ou mais \textit{plays}; cada \textit{play} indica
        em quais \textit{hosts} executar e quais tarefas (\textit{tasks}) aplicar.
  \item \textbf{Task} e \textbf{módulo}: uma \textit{task} é uma ação no servidor (instalar um pacote,
        criar um usuário, copiar um arquivo); ela é executada por um \textbf{módulo}, que é um
        comando pronto do Ansible (\texttt{dnf}, \texttt{user}, \texttt{copy}, \texttt{service}\ldots).
  \item \textbf{Role}: forma de organizar tarefas, variáveis, \textit{templates} e \textit{handlers}
        em um módulo reutilizável.
\end{itemize}

O fluxo de execução é sempre o mesmo: o Ansible lê o inventário, conecta-se via SSH aos
\textit{hosts}, envia e executa os módulos, aplica as configurações e retorna o resultado de cada
tarefa. A Figura~\ref{fig:arquitetura-ansible} mostra como esses elementos foram distribuídos no
laboratório.

\begin{figure}[H]
  \centering
  \caption{Arquitetura do laboratório de Ansible}\label{fig:arquitetura-ansible}
  \begin{tikzpicture}
    \node[caixaescura, text width=4.2cm] (cn) {control-node\\{\normalfont\scriptsize 192.168.1.2 -- Ansible instalado}};
    \node[caixa, text width=4.2cm, below=0.25cm of cn, font=\scriptsize, align=left] (arq)
         {\texttt{inventory/hosts.ini}\\\texttt{playbooks/db.yml} e \texttt{app.yml}\\
          \texttt{roles/} e \texttt{vars/}\\\texttt{\textasciitilde/.ssh/id\_rsa} (chave privada)};
    \begin{scope}[on background layer]
      \node[draw=corprimaria, fill=corprimaria!3, rounded corners=4pt, fit=(cn)(arq), inner sep=6pt] (cnbox) {};
    \end{scope}
    \node[caixalaranja, text width=3.9cm, minimum height=1.35cm, right=3.6cm of cn.north east, anchor=north west] (app)
         {\textbf{app01} -- 192.168.1.3\\{\scriptsize Java 21, Maven e aplicação\\Easy Notes (porta 8080)}};
    \node[caixalaranja, text width=3.9cm, minimum height=1.35cm, below=1.3cm of app] (db)
         {\textbf{db01} -- 192.168.1.4\\{\scriptsize MariaDB (porta 3306)\\banco de dados \texttt{notes}}};
    \draw[seta] ([yshift=0.45cm]cnbox.east) -- node[rotulo, above, sloped]{SSH (22)} node[rotulo, below, sloped]{\texttt{app.yml}} (app.west);
    \draw[seta] ([yshift=-0.45cm]cnbox.east) -- node[rotulo, above, sloped]{SSH (22)} node[rotulo, below, sloped]{\texttt{db.yml}} (db.west);
    \draw[seta, dashed, color=corverde] (app.south) -- node[rotulo, right=2pt]{JDBC (3306)} (db.north);
    \node[font=\scriptsize\bfseries, text=corprimaria, above=1pt of cnbox] {Control Node};
    \node[font=\scriptsize\bfseries, text=cordestaque, above=1pt of app] {Managed Hosts};
  \end{tikzpicture}
  \fonte{Elaborado pelo autor (2026), com base em Muniz (2026).}
\end{figure}

\subsubsection{Preparação do control-node}

O \texttt{control-node} foi criado pelo Vagrant com o script do Código~\ref{cod:provision-cn}. Ele
instala o repositório EPEL (\textit{Extra Packages for Enterprise Linux}), que fornece o pacote do
Ansible para o Rocky Linux, instala o Ansible e acrescenta ao \texttt{/etc/hosts} os nomes das três
máquinas. Assim, o \textit{control-node} consegue alcançar os servidores pelo nome (\texttt{app01} e
\texttt{db01}) em vez do IP. A pasta do projeto no Windows fica disponível na VM em
\texttt{/vagrant}, de modo que os arquivos do Ansible podiam ser editados no VS Code e executados
diretamente no \textit{control-node}.

\codigo{bash}{02-ansible-lab/control-node/provision.sh}{Provisionamento do control-node}{provision-cn}

\subsubsection{Comunicação via SSH}

Com o Ansible instalado, o primeiro teste de conectividade, \cod|ansible -m ping all|, falhou com
\textit{Permission denied}: os servidores ainda não confiavam no \textit{control-node}, porque não
existia autenticação por chave. A solução foi configurar a autenticação SSH por par de chaves
(assimétrica), que é a forma recomendada -- autenticação por senha não é indicada \citep{muniz2026}:

\begin{enumerate}
  \item No \texttt{control-node}, foi gerado o par de chaves com \cod|ssh-keygen|, que cria a chave
        \textbf{privada} (\texttt{\textasciitilde/.ssh/id\_rsa}, que nunca sai da máquina) e a chave
        \textbf{pública} (\texttt{\textasciitilde/.ssh/id\_rsa.pub}, que pode ser compartilhada).
  \item O conteúdo da chave pública (\cod|cat ~/.ssh/id_rsa.pub|) foi colado no
        \texttt{provision.sh} do \texttt{app01} e do \texttt{db01}
        (Código~\ref{cod:provision-app01}), que a adiciona ao arquivo
        \path{/home/vagrant/.ssh/authorized_keys}.
  \item O provisionamento foi reaplicado nas duas máquinas com \cod|vagrant reload --provision|
        (ou \cod|vagrant destroy -f| e \cod|vagrant up|).
  \item A partir daí, o \texttt{control-node} passou a acessar os servidores sem senha, por
        exemplo com \cod|ssh vagrant@app01|, e o \cod|ansible all -m ping| passou a responder
        \texttt{SUCCESS} com \texttt{"ping": "pong"} para os dois \textit{hosts}.
\end{enumerate}

\codigo{bash}{02-ansible-lab/app01/provision.sh}{Provisionamento do app01 com a chave pública do control-node}{provision-app01}

\begin{atencao}[Segurança]
  A chave pública aparece no \texttt{provision.sh}, mas a chave privada (\texttt{id\_rsa}) nunca deve
  ser exibida em prints nem versionada. No repositório, a chave pública real foi substituída pelo
  marcador \texttt{COLE\_AQUI\_A\_CHAVE\_PUBLICA\_DO\_CONTROL\_NODE}.
\end{atencao}

\begin{terminal}[caption={Configuração e teste da comunicação SSH no control-node},label={cod:ssh-ansible}]
vagrant ssh                      # entra no control-node (a partir da pasta control-node)
ansible -m ping all              # 1º teste: falha (Permission denied)
ssh-keygen                       # gera o par de chaves (Enter para aceitar os padrões)
cat ~/.ssh/id_rsa.pub            # exibe a chave PÚBLICA para colar no provision.sh
# ... no Windows, em app01/ e db01/: vagrant reload --provision ...
ssh vagrant@app01                # acesso sem senha (digite exit para voltar)
ansible all -m ping              # app01 | SUCCESS e db01 | SUCCESS ("pong")
\end{terminal}

\subsubsection{Inventário e arquivo de configuração}

No início, o inventário foi configurado no arquivo padrão \texttt{/etc/ansible/hosts}, editado
com o \texttt{vi} (\texttt{i} para inserir, \texttt{ESC} para sair do modo de inserção e
\texttt{:wq} para salvar), com os grupos \texttt{[apps]} (contendo \texttt{app01}) e \texttt{[dbs]}
(contendo \texttt{db01}). Com a evolução do laboratório, o projeto ganhou uma estrutura própria e o
inventário passou a ficar dentro dele, em \texttt{inventory/hosts.ini}
(Código~\ref{cod:inventario}). Os grupos permitem executar comandos e \textit{playbooks} em conjuntos
de servidores (por exemplo, \cod|ansible db -m ping|).

\codigo{ini}{02-ansible-lab/control-node/ansible/inventory/hosts.ini}{Inventário do laboratório}{inventario}

O arquivo de configuração do Ansible (Código~\ref{cod:ansible-cfg}) indica onde estão o inventário
e as \textit{roles} e desativa a verificação interativa da impressão digital (\textit{fingerprint})
das chaves dos \textit{hosts} (\texttt{host\_key\_checking = False}), o que é aceitável em
laboratório. Um problema encontrado foi que o Ansible \textbf{ignorava} esse arquivo quando ele
estava na pasta compartilhada \texttt{/vagrant}, por considerá-la \textit{world writable} (gravável
por qualquer usuário) -- uma proteção de segurança. A solução foi copiá-lo para a \textit{home} do
usuário: \cod|cp /vagrant/ansible/.ansible.cfg /home/vagrant/.ansible.cfg|.

% Conteúdo do arquivo codigos/02-ansible-lab/control-node/ansible/.ansible.cfg
\begin{listagem}[language=ini,caption={[Arquivo de configuração do Ansible (.ansible.cfg)]{Arquivo de configuração do Ansible -- \nolinkurl{02-ansible-lab/control-node/ansible/.ansible.cfg}}},label={cod:ansible-cfg}]
[defaults]
inventory = /vagrant/ansible/inventory/hosts.ini
roles_path = /vagrant/ansible/roles
host_key_checking = False
\end{listagem}

\printtela{ansible-inventario}{Inventário do Ansible no control-node}%
  {Terminal no control-node com \texttt{cat /vagrant/ansible/inventory/hosts.ini}
   (ou \texttt{ansible-inventory -{}-graph}) mostrando os grupos \texttt{[db]} e \texttt{[app]}.}

\printtela{ansible-ssh-keygen}{Geração do par de chaves SSH no control-node}%
  {Saída do \texttt{ssh-keygen} e do \texttt{cat \textasciitilde/.ssh/id\_rsa.pub}.
   NÃO mostrar o conteúdo da chave privada (\texttt{id\_rsa}).}

\printtela{ansible-ping}{Teste de conectividade com \texttt{ansible all -m ping}}%
  {Terminal com \texttt{ansible all -m ping} retornando \texttt{app01 | SUCCESS} e
   \texttt{db01 | SUCCESS} com \texttt{"ping": "pong"}.}

\subsubsection{Estrutura do projeto Ansible}

Para manter a automação padronizada, modular e reutilizável, o projeto foi organizado na estrutura
de diretórios da Figura~\ref{fig:estrutura-ansible}. Cada pasta tem uma responsabilidade:
\texttt{inventory/} define \textbf{quais} máquinas são gerenciadas, \texttt{playbooks/} define
\textbf{o que} será executado e onde, \texttt{vars/} define \textbf{com quais valores}, e
\texttt{roles/} organiza \textbf{como} a configuração é feita, em módulos reutilizáveis
\citep{muniz2026}.

\begin{figure}[H]
  \centering
  \caption{Estrutura de diretórios do projeto Ansible no control-node}\label{fig:estrutura-ansible}
  \begin{minipage}{0.78\linewidth}
    \dirtree{%
      .1 ansible/.
      .2 {.}ansible{.}cfg \DTcomment{configuração (inventário e caminho das roles)}.
      .2 inventory/.
      .3 hosts{.}ini \DTcomment{servidores e grupos}.
      .2 vars/.
      .3 main{.}yml \DTcomment{variáveis globais do projeto}.
      .2 playbooks/.
      .3 db{.}yml \DTcomment{configura o db01}.
      .3 app{.}yml \DTcomment{configura o app01 e a aplicação}.
      .3 templates/ \DTcomment{arquivos Jinja2 usados pelo app{.}yml}.
      .4 application{.}properties.
      .4 etc/systemd/system/notes{.}service.
      .2 roles/.
      .3 configuracao-default-so/ \DTcomment{preparação do sistema}.
      .4 tasks/main{.}yml.
      .4 defaults/ files/ handlers/ templates/ vars/.
      .3 configuracao-mariadb/ \DTcomment{banco de dados}.
      .4 tasks/main{.}yml.
      .4 vars/main{.}yml.
      .4 defaults/ files/ handlers/ templates/.
    }
  \end{minipage}
  \fonte{Elaborado pelo autor (2026), com base em Muniz (2026).}
\end{figure}

Dentro de cada \textit{role}, os nomes das pastas são padronizados e precisam estar exatamente no
plural (\texttt{tasks}, \texttt{handlers}\ldots), porque o Ansible os descobre automaticamente:
ao chamar a \textit{role} no \textit{playbook}, ele executa o arquivo \texttt{tasks/main.yml}. A
Tabela~\ref{tab:pastas-role} resume a função de cada pasta -- ou, na frase usada em aula,
``\textit{tasks} faz, \textit{handlers} reage, \textit{vars} define e \textit{defaults} permite mudar''.

\begin{table}[H]
  \centering
  \caption{Pastas de uma role do Ansible}\label{tab:pastas-role}
  \small
  \renewcommand{\arraystretch}{1.2}
  \begin{tabularx}{\linewidth}{>{\ttfamily}l >{\raggedright\arraybackslash}X}
    \toprule
    \normalfont\bfseries Pasta & \textbf{Função} \\
    \midrule
    tasks/ & Tarefas executadas pela \textit{role} (arquivo \texttt{main.yml}). \\
    handlers/ & Tarefas disparadas somente quando outra tarefa altera algo (\texttt{notify}), como
      reiniciar um serviço. \\
    templates/ & Arquivos dinâmicos (Jinja2), com variáveis substituídas na execução. \\
    files/ & Arquivos estáticos, copiados exatamente como estão (scripts, SQL\ldots). \\
    vars/ & Variáveis específicas da \textit{role} (alta precedência). \\
    defaults/ & Valores padrão (baixa precedência): ``se ninguém informar outro valor, use este''. \\
    \bottomrule
  \end{tabularx}
  \fonte{Elaborado pelo autor (2026), com base em Muniz (2026).}
\end{table}

\subsubsection{Variáveis}

As variáveis do projeto ficam em \texttt{vars/main.yml} (Código~\ref{cod:vars}) e são importadas
pelos \textit{playbooks} com \texttt{vars\_files}. Elas evitam repetir valores nas tarefas: o nome do
banco, o usuário e as senhas são definidos em um único lugar e usados com a sintaxe Jinja2
\cod|"{{ mysql_db_name }}"|. A \textit{role} \path{configuracao-mariadb} também possui o seu
\texttt{vars/main.yml} com os mesmos nomes; como as variáveis de \textit{role} têm precedência maior
que as de \texttt{vars\_files}, os valores precisam ser iguais nos dois arquivos \citep{redhat2026}.

\codigo{yaml}{02-ansible-lab/control-node/ansible/vars/main.yml}{Variáveis globais do projeto}{vars}

\begin{nota}
  Nas aulas, as senhas do laboratório foram gravadas diretamente nesses arquivos. No repositório elas
  foram trocadas por marcadores. Em um ambiente real, o correto é proteger esse arquivo com o
  \textbf{Ansible Vault} (\cod|ansible-vault encrypt vars/main.yml|) e executar os
  \textit{playbooks} com \cod|--ask-vault-pass|.
\end{nota}

\subsubsection{Role configuracao-default-so}

Essa \textit{role} prepara o sistema operacional e é compartilhada pelos dois servidores
(Código~\ref{cod:role-so}): atualiza todos os pacotes, instala o Git, o Python com o PyMySQL (necessário
para os módulos de banco de dados do Ansible), o cliente MariaDB e garante que o \texttt{firewalld}
esteja instalado, habilitado e em execução.

\codigolongo{yaml}{02-ansible-lab/control-node/ansible/roles/configuracao-default-so/tasks/main.yml}{Tarefas da role configuracao-default-so}{role-so}

\subsubsection{Role configuracao-mariadb e Ansible Galaxy}

Para o banco de dados, o primeiro caminho apresentado foi o \textbf{Ansible Galaxy}, repositório
\textit{on-line} de \textit{roles} prontas compartilhadas pela comunidade. Foi instalada a
\textit{role} \texttt{geerlingguy.mysql} e a coleção \texttt{community.mysql}, que fornece os
módulos \texttt{mysql\_db} e \texttt{mysql\_user}:

\begin{terminal}[caption={Comandos do Ansible Galaxy executados no control-node},label={cod:galaxy}]
ansible-galaxy install geerlingguy.mysql                # baixa a role para ~/.ansible/roles/
ansible-galaxy collection install community.mysql        # módulos mysql_db e mysql_user
ansible-galaxy list                                      # lista as roles instaladas
\end{terminal}

Em seguida, a instalação do banco foi implementada em uma \textit{role} própria,
\path{configuracao-mariadb} (Código~\ref{cod:role-mariadb}), o que ajudou a entender cada passo:
instalar o MariaDB e o PyMySQL; habilitar e iniciar o serviço; definir a senha do \texttt{root}
(na primeira execução o acesso é feito pelo \textit{socket} local, ainda sem senha); criar o arquivo
\texttt{/root/.my.cnf} com as credenciais do \texttt{root} e permissão \texttt{0600}; aplicar boas
práticas de segurança (remover usuários anônimos e o banco \texttt{test}); criar o banco
\texttt{notes} e o usuário \texttt{notesapp} com acesso a partir de qualquer \textit{host}
(\texttt{\%}); e liberar a porta 3306 no \textit{firewall}. O \texttt{ignore\_errors} e o arquivo
\texttt{.my.cnf} foram adicionados para que o \textit{playbook} pudesse ser reexecutado sem falhar
depois que o \texttt{root} já tinha senha -- um cuidado ligado à idempotência.

\codigolongo{yaml}{02-ansible-lab/control-node/ansible/roles/configuracao-mariadb/tasks/main.yml}{Tarefas da role configuracao-mariadb}{role-mariadb}

\subsubsection{Playbook db.yml: configuração do servidor de banco de dados}

O \textit{playbook} \texttt{db.yml} (Código~\ref{cod:db-yml}) junta tudo para o \texttt{db01}.
Linha a linha: o \texttt{-{}-{}-} abre o documento YAML; o \texttt{- name} inicia uma \textit{play} (o
hífen indica um item de lista, pois um \textit{playbook} é uma lista de \textit{plays});
\texttt{hosts: db01} define onde executar; \texttt{remote\_user: vagrant} é o usuário da conexão SSH;
\texttt{become: true} executa as tarefas com privilégio de administrador (equivalente ao
\texttt{sudo}); \texttt{vars\_files} importa as variáveis; e \texttt{roles} lista as \textit{roles}
aplicadas, na ordem.

\codigo{yaml}{02-ansible-lab/control-node/ansible/playbooks/db.yml}{Playbook db.yml}{db-yml}

Antes da execução real, o \textit{playbook} foi validado com as opções da
Tabela~\ref{tab:testes-playbook}, todas executadas no \texttt{control-node} a partir da pasta
\texttt{/vagrant/ansible}. Como YAML é sensível à indentação (não aceita tabulação e exige
espaços consistentes, normalmente dois por nível), o \cod|--syntax-check| foi muito útil para achar
erros antes de tocar nos servidores.

\begin{table}[H]
  \centering
  \caption{Comandos usados para testar e executar os playbooks}\label{tab:testes-playbook}
  \small
  \renewcommand{\arraystretch}{1.2}
  \begin{tabularx}{\linewidth}{>{\ttfamily\small}l >{\raggedright\arraybackslash}X}
    \toprule
    \normalfont\bfseries Comando & \textbf{Função} \\
    \midrule
    ansible-playbook playbooks/db.yml -{}-syntax-check & Valida a sintaxe do YAML/\textit{playbook}. \\
    ansible-playbook playbooks/db.yml -{}-list-hosts & Mostra em quais \textit{hosts} o \textit{playbook} vai rodar. \\
    ansible-playbook playbooks/db.yml -{}-list-tasks & Lista as tarefas e \textit{roles} que serão executadas. \\
    ansible-playbook playbooks/db.yml -{}-check & Simula a execução sem aplicar mudanças (\textit{dry run}). \\
    ansible-playbook playbooks/db.yml & Execução real. \\
    ansible-playbook playbooks/db.yml -v & Execução com mais detalhes (\texttt{-vvv} para depuração). \\
    \bottomrule
  \end{tabularx}
  \fonte{Elaborado pelo autor (2026), com base em Muniz (2026).}
\end{table}

Ao final de cada execução, o Ansible exibe o \textbf{PLAY RECAP}, um resumo por \textit{host}:
\texttt{ok} (tarefas que já estavam no estado desejado), \texttt{changed} (tarefas que alteraram
algo), \texttt{unreachable} (\textit{host} inacessível), \texttt{failed} (tarefas com erro),
\texttt{skipped} (puladas) e \texttt{ignored} (erros ignorados com \texttt{ignore\_errors}).

\printtela{ansible-db-playbook}{Execução do playbook db.yml no db01}%
  {Terminal com \texttt{ansible-playbook playbooks/db.yml} mostrando as tarefas das duas
   \textit{roles} e o PLAY RECAP do \texttt{db01} com \texttt{failed=0}.}

\subsubsection{Idempotência}

\textbf{Idempotência} é a capacidade de executar a mesma tarefa várias vezes e sempre obter o mesmo
resultado final: se o servidor já está no estado desejado, nada é alterado \citep{muniz2026}. Na
prática, isso ficou visível ao executar o \texttt{db.yml} duas vezes: na primeira execução várias
tarefas apareceram como \texttt{changed} (o MariaDB foi instalado, o banco foi criado\ldots); na
segunda, as mesmas tarefas apareceram como \texttt{ok}, com \texttt{changed=0} ou próximo de zero,
porque o banco, o usuário e as regras de \textit{firewall} já existiam. Essa propriedade evita
erros em reexecuções, garante a consistência do ambiente e é o que permite usar o Ansible com
segurança dentro de \textit{pipelines} de CI/CD. Vale notar que nem todo módulo é idempotente por
natureza: tarefas com \texttt{command} ou \texttt{shell} (como o \cod|mvn package| do
\texttt{app.yml}) sempre aparecem como \texttt{changed}, a menos que se use \texttt{changed\_when}
ou \texttt{creates} -- por isso a tarefa que apenas localiza o JAR recebeu
\texttt{changed\_when: false}.

\printtela{ansible-idempotencia}{Segunda execução do db.yml demonstrando a idempotência}%
  {Terminal com a segunda execução do \texttt{ansible-playbook playbooks/db.yml}, mostrando as
   tarefas em verde (\texttt{ok}) e o PLAY RECAP com \texttt{changed=0} (ou o menor valor obtido).}

\subsubsection{Playbook app.yml: servidor de aplicação e deploy do Easy Notes}

O segundo \textit{playbook} configura o \texttt{app01} e faz o \textit{deploy} da aplicação
\textbf{Easy Notes}, uma API REST de notas desenvolvida em Java com Spring Boot, que grava os dados
no MariaDB do \texttt{db01} \citep{singh2021}. O projeto (pasta \texttt{Easy-Notes-Java},
disponibilizada no Blackboard) foi salvo dentro da pasta do \texttt{control-node}. O fluxo do
\textit{playbook} é:

\begin{tcolorbox}[enhanced, colback=corprimaria!4, colframe=corprimaria!40, boxrule=0.5pt,
    arc=2pt, left=8pt, right=8pt, top=5pt, bottom=5pt, fontupper=\small\raggedright]
  configurar o \texttt{/etc/hosts} $\rightarrow$ criar o usuário \texttt{app} $\rightarrow$
  instalar o Java 21 $\rightarrow$ definir \texttt{java} e \texttt{javac} padrão $\rightarrow$
  instalar o Maven $\rightarrow$ criar \texttt{/opt/notes} $\rightarrow$ copiar o projeto
  $\rightarrow$ gerar o \texttt{application.properties} $\rightarrow$ compilar com o Maven
  $\rightarrow$ localizar o JAR $\rightarrow$ criar o serviço \texttt{notes.service}
  $\rightarrow$ liberar a porta 8080 $\rightarrow$ iniciar a aplicação.
\end{tcolorbox}

\codigolongo{yaml}{02-ansible-lab/control-node/ansible/playbooks/app.yml}{Playbook app.yml}{app-yml}

Os principais pontos do \texttt{app.yml} são:

\begin{itemize}
  \item \textbf{\texttt{lineinfile} com \texttt{loop}}: garante que cada linha da lista exista no
        \texttt{/etc/hosts}; o \cod|{{ item }}| representa o item atual do \texttt{loop}.
  \item \textbf{\texttt{user}}: cria o usuário \texttt{app}, que executa a aplicação. Assim, a
        aplicação não roda como \texttt{root}, o que é uma boa prática de segurança.
  \item \textbf{\texttt{dnf}} e \textbf{\texttt{alternatives}}: instalam o Java 21 e o definem como
        versão padrão de \texttt{java} e \texttt{javac} (\texttt{priority: 2100}). O caminho exato
        do Java depende da versão instalada e pode ser conferido com \cod|ls /usr/lib/jvm/|.
  \item \textbf{\texttt{file}} e \textbf{\texttt{copy}}: criam \texttt{/opt/notes} com dono
        \texttt{app} e permissão \texttt{0755} e copiam o projeto do \textit{control-node} para o
        \texttt{app01}.
  \item \textbf{\texttt{template}}: gera o \texttt{application.properties} a partir do
        \textit{template} Jinja2 (Código~\ref{cod:app-properties}), substituindo
        \cod|{{ dbhost }}|, \cod|{{ dbname }}|, \cod|{{ dbusername }}| e \cod|{{ dbpassword }}|.
  \item \textbf{\texttt{command}} com \texttt{become\_user: app} e \texttt{environment}: compila a
        aplicação com \cod|mvn clean package| como usuário \texttt{app}, usando o Java 21
        (\texttt{JAVA\_HOME}). Executar como \texttt{app} evita erros de permissão e o erro
        \textit{detected dubious ownership} do Git.
  \item \textbf{\texttt{shell}} com \texttt{register} e \textbf{\texttt{debug}}: localizam o arquivo
        JAR gerado, guardam o caminho na variável \texttt{app\_jar} e o exibem na tela.
  \item \textbf{\texttt{template}} com \texttt{notify}: cria o serviço \texttt{systemd}
        (Código~\ref{cod:notes-service}) e, se o arquivo mudar, aciona os \textit{handlers}.
  \item \textbf{\texttt{firewalld}} e \textbf{\texttt{service}}: liberam a porta 8080 e iniciam e
        habilitam o serviço \texttt{notes} no \textit{boot}.
  \item \textbf{\textit{Handlers}}: \texttt{reload daemon} (recarrega o \texttt{systemd} para ler o
        novo arquivo de serviço) e \texttt{restart notes} (reinicia a aplicação). Eles só executam
        quando notificados e no final da \textit{play}, evitando reinicializações desnecessárias.
\end{itemize}

\codigo{ini}{02-ansible-lab/control-node/ansible/playbooks/templates/application.properties}{Template Jinja2 do application.properties}{app-properties}

O arquivo \texttt{notes.service} (Código~\ref{cod:notes-service}) transforma a aplicação Java em um
serviço do Linux: a seção \texttt{[Unit]} descreve o serviço e indica que ele deve subir depois da
rede (\texttt{After=network.target}); a seção \texttt{[Service]} define o usuário, o diretório de
trabalho, o comando de inicialização -- em que \cod|{{ app_jar.stdout }}| é substituído pelo caminho
do JAR encontrado (por exemplo, \texttt{/opt/notes/target/easy-notes-1.0.0.jar}) -- e
\texttt{Restart=always}, que reinicia a aplicação se ela cair; e a seção \texttt{[Install]} permite
iniciá-la automaticamente no \textit{boot}. Depois disso, a aplicação pode ser controlada com
\cod|systemctl start|, \cod|stop|, \cod|restart| e \cod|status notes|.

\codigo{ini}{02-ansible-lab/control-node/ansible/playbooks/templates/etc/systemd/system/notes.service}{Template do serviço systemd da aplicação}{notes-service}

\subsubsection{Validação da aplicação}

Com os dois \textit{playbooks} executados, a aplicação foi testada com o \texttt{curl}, que permite
``conversar'' com a API pelo terminal: \texttt{-X} define o método HTTP, \texttt{-H} adiciona
cabeçalhos (como \texttt{Content-Type: application/json}) e \texttt{-d} envia o corpo da requisição
(Código~\ref{cod:testes-api}). Por fim, foi conferido no \texttt{db01} que a nota criada pela API
estava gravada na tabela \texttt{notes}, o que comprova a integração entre os dois servidores.

\begin{terminal}[caption={Testes da aplicação e do banco de dados},label={cod:testes-api}]
ansible-playbook playbooks/app.yml          # no control-node, a partir de /vagrant/ansible
ssh vagrant@app01                           # acessa o servidor de aplicação
sudo systemctl status notes                 # serviço ativo (running)
curl http://localhost:8080/api/notes        # lista as notas (JSON)
curl -X POST http://localhost:8080/api/notes \
  -H "Content-Type: application/json" \
  -d '{"title": "Primeira nota", "content": "Aplicação configurada com Ansible"}'

ssh vagrant@db01                            # no servidor de banco de dados
mysql -u notesapp -p notes -e "SELECT id, title FROM notes;"
\end{terminal}

\printtela{ansible-app-playbook}{Execução do playbook app.yml no app01}%
  {Terminal com \texttt{ansible-playbook playbooks/app.yml}: tarefas executadas, o \texttt{debug}
   com o caminho do JAR, os \textit{handlers} (\texttt{RUNNING HANDLER}) e o PLAY RECAP com
   \texttt{failed=0}.}

\printtela{ansible-notes-status}{Serviço da aplicação ativo no app01}%
  {Terminal no app01 com \texttt{sudo systemctl status notes} mostrando \texttt{active (running)}.}

\printtela{ansible-api-curl}{Teste da API REST com curl}%
  {Terminal com o \texttt{curl -X POST} criando uma nota e o \texttt{curl} GET em
   \texttt{/api/notes} retornando o JSON com a nota criada.}

\printtela{ansible-mariadb}{Nota gravada no banco de dados MariaDB do db01}%
  {Terminal no db01 com \texttt{sudo systemctl status mariadb} e/ou o \texttt{SELECT} na tabela
   \texttt{notes} mostrando a nota criada pela API.}

\printtela{ansible-arquivos-yml}{Arquivos YAML do projeto Ansible no editor}%
  {VS Code com a árvore do projeto \texttt{ansible/} aberta e um dos arquivos \texttt{.yml}
   (por exemplo, \texttt{app.yml} ou \texttt{db.yml}) visível.}

\subsubsection{Módulos utilizados}

\begin{table}[H]
  \centering
  \caption{Módulos do Ansible utilizados no laboratório}\label{tab:modulos}
  \small
  \renewcommand{\arraystretch}{1.18}
  \begin{tabularx}{\linewidth}{>{\ttfamily}l >{\raggedright\arraybackslash}X}
    \toprule
    \normalfont\bfseries Módulo & \textbf{Função no laboratório} \\
    \midrule
    dnf & Instala, atualiza ou remove pacotes no Rocky Linux (\texttt{present}, \texttt{latest}, \texttt{absent}). \\
    service / systemd & Inicia, para, reinicia e habilita serviços (\texttt{started}, \texttt{restarted}, \texttt{enabled}). \\
    lineinfile & Garante que uma linha exista em um arquivo (\texttt{/etc/hosts}). \\
    user & Cria e configura usuários (usuário \texttt{app}). \\
    file & Cria diretórios e ajusta dono, grupo e permissões. \\
    copy & Copia arquivos e pastas do \textit{control-node} para o \textit{host}. \\
    template & Copia arquivos processando variáveis Jinja2. \\
    command / shell & Executa comandos (\texttt{shell} permite \textit{pipes}, como \texttt{|} e \texttt{grep}). \\
    debug & Exibe valores de variáveis na saída do \textit{playbook}. \\
    alternatives & Define a versão padrão de um programa (\texttt{java}, \texttt{javac}). \\
    firewalld & Abre portas e serviços no \textit{firewall}. \\
    mysql\_db / mysql\_user & Criam e removem bancos e usuários no MariaDB (coleção \texttt{community.mysql}). \\
    \bottomrule
  \end{tabularx}
  \fonte{Elaborado pelo autor (2026), com base em Muniz (2026).}
\end{table}

\subsubsection{Diagnóstico de problemas}

Quando algo não funcionava, foi seguida a sequência de verificação apresentada em aula, do mais
básico ao mais específico \citep{muniz2026}: (1) a máquina responde? (\cod|ping app01|); (2) a porta
está aberta? (\cod!sudo ss -tulpn | grep 8080! ou \cod|nc -zv app01 8080|); (3) o
\textit{firewall} está liberado? (\cod|sudo firewall-cmd --list-all|); (4) o serviço está rodando?
(\cod|systemctl status notes| e \cod|journalctl -u notes|); (5) a API responde localmente?
(\cod|curl http://localhost:8080/api/notes|); (6) responde remotamente? (a partir do
\textit{control-node}); e (7) o JSON enviado está correto? Essa abordagem em camadas evitou
``chutes'' e ajudou a entender a relação entre rede, \textit{firewall}, serviço e aplicação.

\subsubsection{Como o Ansible automatiza configurações que antes eram manuais}

Sem o Ansible, preparar o ambiente exigiria entrar em cada servidor por SSH e repetir dezenas de
comandos: atualizar o sistema, instalar pacotes, configurar o MariaDB, criar banco e usuário, abrir
portas, instalar Java e Maven, copiar e compilar a aplicação, escrever o arquivo de serviço e
iniciá-lo. Cada um desses passos é uma oportunidade de erro e, se um servidor precisasse ser
recriado, tudo teria de ser refeito. Com o Ansible, essas etapas viraram código: dois
\textit{playbooks} e duas \textit{roles} versionados que configuram os servidores com um único
comando, sempre da mesma forma, e que podem ser reexecutados com segurança graças à idempotência.
Além disso, a \textit{role} \path{configuracao-default-so} foi reaproveitada nos dois servidores,
mostrando na prática o ganho de padronização e reutilização que a Infraestrutura como Código
proporciona.

% =====================================================================
%  2.4 Atividade com Docker
% =====================================================================
\subsection{Atividade com Docker}\label{sec:docker}

O \textbf{Docker} é uma plataforma de conteinerização: empacota a aplicação com todas as suas
dependências em uma \textbf{imagem} (gerada a partir de um \texttt{Dockerfile}) e a executa como um
\textbf{container}, um processo isolado que compartilha o \textit{kernel} do host. Por isso, um
container é muito mais leve que uma VM (megabytes em vez de gigabytes, segundos em vez de minutos para
iniciar) e roda igual em qualquer máquina \citep{docker2026,muniz2026}.

Na prática, foi criada a VM \texttt{docker-notes} (\texttt{192.168.1.5}), com o Docker instalado pelo
script de provisionamento. O JAR da aplicação, já compilado pelo Ansible no \texttt{app01}, foi copiado
para essa VM (com \cod|vagrant ssh-config| e \texttt{scp} e depois \cod|vagrant upload|), junto com o
\texttt{application.properties}, e a imagem foi construída com o \texttt{Dockerfile} do
Código~\ref{cod:dockerfile}: \texttt{FROM} define a imagem base (um Java enxuto sobre Alpine Linux),
\texttt{RUN} e \texttt{USER} criam e usam um usuário sem privilégios, \texttt{COPY} copia o JAR e a
configuração, \texttt{EXPOSE} documenta a porta e \texttt{ENTRYPOINT} define o comando de
inicialização. A versão do Java da imagem precisa ser compatível com a usada para compilar o JAR
(Java 21); com uma versão menor, o container falha com \texttt{UnsupportedClassVersionError}.

\codigo{docker}{03-docker-lab/Dockerfile}{Dockerfile da aplicação Easy Notes}{dockerfile}

A imagem foi construída e executada com os comandos do Código~\ref{cod:docker-comandos}. A opção
\path{--add-host} permite que o container encontre o banco de dados do \texttt{db01}.

\begin{terminal}[caption={Construção e execução do container},label={cod:docker-comandos}]
docker build -t devops/notes .                             # constrói a imagem
docker run -d --name notes -p 8080:8080 \
  --add-host db01:192.168.1.4 devops/notes                 # inicia o container
docker ps                                                  # container em execução
docker logs -f notes                                       # log da aplicação
\end{terminal}

Também foi visto o
\textit{multi-stage build}, em que um primeiro estágio compila o projeto com Maven e o segundo leva
para a imagem final só o JAR, deixando-a menor e mais segura (arquivo \texttt{Dockerfile.multistage}
no repositório).

\printtela{docker-container}{Imagem construída e container da aplicação em execução}%
  {Terminal na VM docker-notes com \texttt{docker images}, \texttt{docker ps} (container
   \texttt{notes} com status \textit{Up}) e o \texttt{curl} em \texttt{/api/notes} respondendo.}

% =====================================================================
%  2.5 Atividade com SonarQube
% =====================================================================
\subsection{Atividade com SonarQube}\label{sec:sonarqube}

\subsubsection{O que é o SonarQube e para que serve a análise estática}

O \textbf{SonarQube} é uma ferramenta de \textbf{análise estática de código}: ele examina o
código-fonte \textbf{sem executar a aplicação} e aponta problemas, exibindo os resultados em
\textit{dashboards} \citep{sonar2026}. A análise estática serve para encontrar cedo -- quando a
correção é mais barata -- \textbf{\textit{bugs}} (erros que causam comportamento incorreto),
\textbf{vulnerabilidades} (falhas exploráveis por um atacante), \textbf{\textit{security hotspots}}
(trechos sensíveis que precisam de revisão), \textbf{\textit{code smells}} (problemas de manutenção),
além de duplicações e baixa cobertura de testes. Um exemplo clássico é uma senha escrita direto no
código. No DevOps, o SonarQube entra no \textit{pipeline} como um ``portão de qualidade''
(\textit{Quality Gate}): se o código não atende aos critérios, a entrega é bloqueada \citep{muniz2026}.

\subsubsection{Criação do ambiente e acesso ao painel}

O SonarQube foi instalado em uma VM própria (\texttt{sonarqube}, \texttt{192.168.1.6}, 4 GB de
memória, porta 9000 redirecionada para o Windows). O script de provisionamento automatizou a
instalação: Java 21, criação do usuário \texttt{sonar} (o Elasticsearch interno não roda como
\texttt{root}), ajuste do \textit{kernel} exigido pelo Elasticsearch (\texttt{vm.max\_map\_count}),
download do SonarQube para \texttt{/opt/sonarqube}, criação do serviço \texttt{sonar} no
\texttt{systemd}, liberação da porta 9000 no \textit{firewall} e instalação do \textbf{SonarScanner
CLI}. O painel foi acessado pelo navegador em \url{http://127.0.0.1:9000}, com o usuário
\texttt{admin}/\texttt{admin}, cuja senha o SonarQube obriga a trocar no primeiro acesso.

\subsubsection{Criação do projeto e execução do sonar-scanner}

No painel, o projeto foi criado em \textit{Create Project} $\rightarrow$ \textit{Local project}, com
o nome \textit{Easy Notes} e a chave \texttt{easy-notes}; como método de análise foi escolhido
\textit{Locally}, e o SonarQube gerou um \textbf{\textit{token}} -- uma credencial que o
\textit{scanner} usa para enviar os resultados. Para projetos Java, o \textit{scanner} precisa também
das classes compiladas, por isso o projeto foi clonado e compilado com o Maven antes da análise
(Código~\ref{cod:sonar-scanner}).

\begin{terminal}[caption={Compilação do projeto e execução do sonar-scanner},label={cod:sonar-scanner}]
git clone https://github.com/callicoder/spring-boot-mysql-rest-api-tutorial.git easy-notes
cd easy-notes && mvn -B clean package -DskipTests     # gera target/classes
export SONAR_TOKEN=************                       # token gerado no painel (oculto)
sonar-scanner \
  -Dsonar.projectKey=easy-notes \
  -Dsonar.sources=src/main/java \
  -Dsonar.java.binaries=target/classes \
  -Dsonar.host.url=http://localhost:9000 \
  -Dsonar.token=$SONAR_TOKEN
\end{terminal}

Os parâmetros indicam a chave do projeto (\texttt{projectKey}), onde está o código
(\texttt{sources}), as classes compiladas (\texttt{java.binaries}), o endereço do servidor
(\texttt{host.url}) e a credencial (\texttt{token}; em versões antigas, \texttt{sonar.login}). Ao
final, o \textit{scanner} exibe \texttt{ANALYSIS SUCCESSFUL} e o resultado aparece no
\textit{dashboard}.

\begin{atencao}[Token]
  O \textit{token} funciona como uma senha: ele foi guardado em uma variável de ambiente, não aparece
  em arquivos versionados e foi ocultado nos prints.
\end{atencao}

\subsubsection{Resultados da análise}

% CONFIRA: preencha a tabela com os números do SEU dashboard (os mesmos do print).
\begin{table}[H]
  \centering
  \caption{Resultado da análise do projeto Easy Notes no SonarQube}\label{tab:resultado-sonar}
  \small
  \renewcommand{\arraystretch}{1.2}
  \begin{tabular}{l c}
    \toprule
    \textbf{Métrica} & \textbf{Resultado} \\
    \midrule
    Quality Gate & \preencher{Passed / Failed} \\
    \textit{Bugs} / Confiabilidade & \preencher{qtd. e nota} \\
    Vulnerabilidades / Segurança & \preencher{qtd. e nota} \\
    \textit{Security Hotspots} & \preencher{qtd.} \\
    \textit{Code Smells} / Manutenibilidade & \preencher{qtd. e nota} \\
    Duplicações e cobertura & \preencher{\% e \%} \\
    \bottomrule
  \end{tabular}
  \fonte{Elaborado pelo autor (2026), a partir do \textit{dashboard} do SonarQube.}
\end{table}

% CONFIRA: ajuste a lista abaixo com os apontamentos que aparecerem na aba Issues do SEU projeto.
Na aba \textit{Issues}, os apontamentos foram principalmente de manutenção e segurança, como: a
\textbf{injeção de dependência por campo} (\texttt{@Autowired} no atributo do
\texttt{NoteController}, quando o recomendado é injetar pelo construtor); a \textbf{entidade JPA
recebida diretamente na API} (\texttt{@RequestBody Note}), risco de \textit{mass assignment} que se
resolve com um objeto de transferência (DTO); o retorno com tipo genérico curinga
(\texttt{ResponseEntity<?>}); e atributos não \texttt{final} na classe de exceção.

\printtela{sonar-token}{Geração do token de análise (valor ocultado)}%
  {Tela de geração do \textit{token} com o valor \textbf{borrado/tarjado}.}

\printtela{sonar-scanner-exec}{Execução do sonar-scanner}%
  {Final da saída do \texttt{sonar-scanner} com \texttt{ANALYSIS SUCCESSFUL} e
   \texttt{EXECUTION SUCCESS} (sem o \textit{token} visível).}

\printtela{sonar-dashboard}{Dashboard do projeto no SonarQube}%
  {Tela \textit{Overview} do projeto com o Quality Gate e as métricas.}

\printtela{sonar-issues}{Problemas apontados na aba Issues}%
  {Aba \textit{Issues} do projeto listando os apontamentos.}

\subsubsection{Como o SonarQube melhora a qualidade e a segurança antes do deploy}

O SonarQube funciona como um fiscal automático: cada versão do código é analisada de forma objetiva,
e os problemas aparecem logo após o \textit{commit}, quando ainda são baratos de corrigir. As regras
aplicadas a cada linguagem formam o \textbf{Quality Profile}, e as condições para aprovar o código
(por exemplo, nenhum problema grave novo e duplicação baixa) formam o \textbf{Quality Gate}, que no
\textit{pipeline} impede que código reprovado chegue ao \textit{deploy}. Ele não substitui testes nem
revisão humana -- não entende a regra de negócio --, mas complementa o processo de qualidade.

% =====================================================================
%  2.6 Atividade com Jenkins
% =====================================================================
\subsection{Atividade com Jenkins}\label{sec:jenkins}

\subsubsection{O que é o Jenkins}

O \textbf{Jenkins} é um servidor de automação \textit{open source}, escrito em Java, usado para
construir \textit{pipelines} de integração e entrega contínuas. Ele executa \textit{jobs} -- conjuntos
de passos como clonar um repositório, compilar, testar e publicar -- de forma manual, agendada ou
disparada por eventos, como um novo \textit{commit}. Por meio de milhares de \textit{plugins}, integra-se
a praticamente todas as ferramentas do ecossistema DevOps (Git, Maven, Docker, SonarQube, Kubernetes),
e permite descrever o \textit{pipeline} como código, em um arquivo chamado \texttt{Jenkinsfile}
\citep{jenkins2026}.

\subsubsection{CI/CD: integração contínua, entrega contínua e deploy contínuo}

\textbf{CI/CD} é o conjunto de práticas que automatiza o caminho do código desde o \textit{commit}
até a produção \citep{humble2014}. As três práticas são complementares e diferem principalmente em
até onde a automação vai (Tabela~\ref{tab:cicd}):

\begin{itemize}
  \item \textbf{Integração Contínua} (\textit{Continuous Integration}): os desenvolvedores integram o
        código ao repositório principal com frequência, e cada integração dispara automaticamente o
        \textit{build} e os testes. Assim, erros de integração são descobertos em minutos, e não
        semanas depois.
  \item \textbf{Entrega Contínua} (\textit{Continuous Delivery}): estende a CI. Toda versão que passa
        pelo \textit{pipeline} fica \textbf{pronta} para ser liberada, mas a ida para produção
        depende de uma \textbf{aprovação manual} (um clique, uma decisão de negócio).
  \item \textbf{Deploy Contínuo} (\textit{Continuous Deployment}): vai um passo além: toda alteração
        aprovada em todas as etapas é implantada \textbf{automaticamente} em produção, sem
        intervenção humana.
\end{itemize}

\begin{table}[H]
  \centering
  \caption{Diferenças entre integração contínua, entrega contínua e deploy contínuo}\label{tab:cicd}
  \small
  \renewcommand{\arraystretch}{1.2}
  \begin{tabularx}{\linewidth}{l >{\raggedright\arraybackslash}X >{\raggedright\arraybackslash}X}
    \toprule
    \textbf{Prática} & \textbf{O que é automatizado} & \textbf{Ida para produção} \\
    \midrule
    Integração contínua & \textit{Build} e testes a cada \textit{commit} & Não se aplica \\
    Entrega contínua & \textit{Build}, testes, análise e empacotamento; artefato sempre pronto & Manual (aprovação) \\
    \textit{Deploy} contínuo & Todo o caminho, inclusive o \textit{deploy} & Automática \\
    \bottomrule
  \end{tabularx}
  \fonte{Elaborado pelo autor (2026), com base em Humble e Farley (2014).}
\end{table}

\subsubsection{Como o Jenkins foi instalado e acessado}

Seguindo o mesmo padrão das demais ferramentas, o Jenkins foi instalado em uma VM própria
(\texttt{jenkins}, \texttt{192.168.1.7}, 4 GB de memória), com a interface web redirecionada para a
porta 8081 do Windows -- a 8080 já era usada pelo \texttt{app01}. O script de provisionamento
(Código~\ref{cod:provision-jenkins}) instala o Java 21, exigido pelas versões atuais do Jenkins, o
Java 17 e o Maven para compilar a aplicação, o Jenkins LTS a partir do repositório oficial, o Docker
(com o usuário \texttt{jenkins} no grupo \texttt{docker}, para que o \textit{pipeline} possa gerar
imagens) e o SonarScanner. Também adiciona ao \texttt{/etc/hosts} os nomes \texttt{app01},
\texttt{db01} e \texttt{sonarqube}, para que o Jenkins alcance as outras VMs pelo nome.

\codigolongo{bash}{05-jenkins/provision.sh}{Provisionamento da VM do Jenkins}{provision-jenkins}

Depois do \cod|vagrant up|, o Jenkins foi acessado em \url{http://localhost:8081} e configurado pelo
assistente inicial:

\begin{enumerate}
  \item desbloqueio (\textit{Unlock Jenkins}) com a senha inicial, lida na VM com
        \cod|sudo cat /var/lib/jenkins/secrets/initialAdminPassword|;
  \item instalação dos \textit{plugins} sugeridos (\textit{Install suggested plugins}), que incluem
        Pipeline, Git e Credentials Binding;
  \item criação do primeiro usuário administrador e confirmação da URL do Jenkins.
\end{enumerate}

Para não escrever segredos nos \textit{pipelines}, as credenciais foram cadastradas em
\textit{Manage Jenkins} $\rightarrow$ \textit{Credentials}: o \textit{token} do SonarQube como
\textit{Secret text} (ID \texttt{sonar-token}) e o usuário e a senha do banco da aplicação como
\textit{Username with password} (ID \texttt{db-notes}). No \textit{pipeline}, elas são lidas com
\texttt{withCredentials}, e o Jenkins ainda mascara os valores nos logs.

\printtela{jenkins-inicial}{Tela inicial (dashboard) do Jenkins}%
  {Navegador em \texttt{http://localhost:8081} mostrando o \textit{dashboard} do Jenkins com os
   \textit{jobs} criados. (Se fizer print do desbloqueio, tarje a senha inicial.)}

\subsubsection{Como foi criado o primeiro job}

O primeiro \textit{job} foi do tipo \textit{Freestyle project}, configurado pela interface:
\textit{New Item} $\rightarrow$ nome \texttt{primeiro-job} $\rightarrow$ \textit{Freestyle project}
$\rightarrow$ em \textit{Build Steps}, \textit{Add build step} $\rightarrow$ \textit{Execute shell},
com o script do Código~\ref{cod:primeiro-job}. Ao clicar em \textit{Build Now}, o Jenkins executou o
script na própria VM e registrou a execução no histórico (\textit{Build History}); o
\textit{Console Output} mostrou cada comando e terminou com \texttt{Finished: SUCCESS}. As variáveis
\texttt{JOB\_NAME} e \texttt{BUILD\_NUMBER} são fornecidas automaticamente pelo Jenkins.

\codigo{bash}{05-jenkins/primeiro-job.sh}{Script do passo Execute shell do primeiro job}{primeiro-job}

\printtela{jenkins-primeiro-job}{Configuração do primeiro job (Freestyle)}%
  {Tela de configuração do \texttt{primeiro-job} mostrando o passo \textit{Execute shell} com o
   script, ou a página do \textit{job} com o \textit{Build History}.}

\printtela{jenkins-job-console}{Console Output do primeiro job}%
  {\textit{Console Output} do \textit{build} \#1 do \texttt{primeiro-job} com a saída dos comandos e
   \texttt{Finished: SUCCESS}.}

\subsubsection{Como foi criado um pipeline}

Em seguida foi criado um \textit{job} do tipo \textit{Pipeline} (\textit{New Item} $\rightarrow$
\texttt{pipeline-basico} $\rightarrow$ \textit{Pipeline}), e, em \textit{Definition}, foi escolhida
a opção \textit{Pipeline script} com o \texttt{Jenkinsfile} do Código~\ref{cod:jenkinsfile-basico}.
Esse \textit{pipeline} usa a sintaxe \textbf{declarativa}: o bloco \texttt{pipeline} envolve toda a
definição; \texttt{agent any} indica que ele pode rodar em qualquer executor disponível;
\texttt{stages} agrupa as etapas, executadas em sequência; cada \texttt{stage} é uma etapa nomeada,
que aparece na visualização dos estágios; e \texttt{steps} contém os comandos de cada etapa -- aqui,
apenas \texttt{echo}, que imprime uma mensagem no log. Mesmo simples, ele mostra a estrutura de
qualquer \textit{pipeline} real: \textit{Build}, \textit{Test} e \textit{Deploy}.

\codigo{groovy}{05-jenkins/Jenkinsfile.basico}{Jenkinsfile do pipeline básico}{jenkinsfile-basico}

\subsubsection{Pipeline integrado: Git, Maven, SonarQube e Docker}

O passo seguinte foi transformar o exemplo em um \textit{pipeline} real para a aplicação Easy Notes,
integrando as ferramentas estudadas (Figura~\ref{fig:pipeline} e Código~\ref{cod:jenkinsfile}).
O \textit{job} \texttt{easy-notes-ci} foi criado do mesmo modo, colando o \texttt{Jenkinsfile} em
\textit{Pipeline script}; a alternativa \textit{Pipeline script from SCM} faz o Jenkins ler o
\texttt{Jenkinsfile} direto do repositório Git, o que deixa o próprio \textit{pipeline} versionado
(\textit{pipeline as code}).

\begin{figure}[H]
  \centering
  \caption{Estágios do pipeline de CI/CD da aplicação Easy Notes}\label{fig:pipeline}
  \begin{tikzpicture}[node distance=0.55cm]
    \tikzset{est/.style={caixa, text width=2.35cm, minimum height=1.15cm}}
    \node[caixalaranja, text width=2.35cm, minimum height=1.15cm] (c) {\textbf{Commit / push}\\{\scriptsize repositório Git}};
    \node[est, right=of c] (s1) {\textbf{Checkout}\\{\scriptsize \texttt{git}}};
    \node[est, right=of s1] (s2) {\textbf{Build}\\{\scriptsize \texttt{mvn package}}};
    \node[est, right=of s2] (s3) {\textbf{Testes}\\{\scriptsize \texttt{mvn test} (db01)}};
    \node[est, below=1.15cm of s3] (s4) {\textbf{SonarQube}\\{\scriptsize análise e \textit{quality gate}}};
    \node[est, left=of s4] (s5) {\textbf{Imagem Docker}\\{\scriptsize \texttt{docker build}}};
    \node[est, left=of s5] (s6) {\textbf{Deploy}\\{\scriptsize \texttt{docker run}}};
    \node[caixa, draw=corverde, fill=corverde!8, text width=2.35cm, minimum height=1.15cm, left=of s6] (s7)
         {\textbf{Teste de fumaça}\\{\scriptsize \texttt{curl /api/notes}}};
    \draw[seta] (c) -- node[rotulo, above=1pt]{\scriptsize poll} (s1);
    \draw[seta] (s1) -- (s2);
    \draw[seta] (s2) -- (s3);
    \draw[seta] (s3) -- (s4);
    \draw[seta] (s4) -- node[rotulo, above=1pt]{aprovado} (s5);
    \draw[seta] (s5) -- (s6);
    \draw[seta] (s6) -- (s7);
    \node[caixa, draw=corvermelha, fill=corvermelha!7, text width=3.7cm, below=0.85cm of s4] (f)
         {\textbf{Pipeline interrompido}\\{\scriptsize o time corrige e envia um novo \textit{commit}}};
    \draw[seta, color=corvermelha] (s4) -- node[rotulo, right=2pt]{reprovado} (f);
    \draw[seta, dashed, color=corvermelha] (f.west) -| ($(s7.west)+(-0.45,0)$) |- (c.west);
    \node[rotulo, rotate=90, anchor=south] at ($(s7.west)+(-0.45,-0.9)$) {correção};
  \end{tikzpicture}
  \fonte{Elaborado pelo autor (2026).}
\end{figure}

\codigolongo{groovy}{05-jenkins/Jenkinsfile}{Jenkinsfile do pipeline integrado da aplicação Easy Notes}{jenkinsfile}

Cada parte do \texttt{Jenkinsfile} cumpre um papel na integração:

\begin{itemize}
  \item \textbf{\texttt{environment}}: centraliza endereços e nomes (repositório, servidor do
        SonarQube, URL do banco, nome da imagem) e define o \texttt{JAVA\_HOME} usado pelo Maven.
  \item \textbf{\texttt{triggers}} com \texttt{pollSCM}: o Jenkins consulta o repositório a cada cinco
        minutos e dispara o \textit{pipeline} quando encontra um \textit{commit} novo -- integração
        com o \textbf{Git}. (Um \textit{webhook} do GitHub seria mais imediato, mas exigiria que o
        GitHub alcançasse a VM, o que não ocorre em uma rede privada de laboratório.)
  \item \textbf{Checkout}: o passo \texttt{git} clona o repositório da aplicação no
        \textit{workspace} do \textit{job}.
  \item \textbf{Build}: compila e empacota a aplicação com o Maven (\cod|mvn clean package|).
  \item \textbf{Testes}: executa os testes; o teste da aplicação sobe o contexto do Spring e conecta-se
        ao MariaDB do \texttt{db01}, configurado pelo Ansible. As credenciais chegam pelas variáveis
        \texttt{SPRING\_DATASOURCE\_*}, sem ficarem escritas no código.
  \item \textbf{Análise SonarQube}: executa o \texttt{sonar-scanner} com o \textit{token} vindo das
        credenciais do Jenkins -- integração com o \textbf{SonarQube}. Com
        \texttt{sonar.qualitygate.wait=true}, o \textit{scanner} aguarda o resultado do Quality Gate e
        o estágio falha se o código for reprovado, interrompendo o \textit{pipeline}.
  \item \textbf{Imagem Docker}: gera o \texttt{Dockerfile} e constrói a imagem
        \texttt{devops/notes}, marcada com o número do \textit{build} -- integração com o
        \textbf{Docker}.
  \item \textbf{Deploy} e \textbf{Teste de fumaça}: substituem o container anterior pelo novo e
        confirmam, com \texttt{curl}, que a API responde; o \texttt{retry} repete a verificação
        enquanto a aplicação inicializa.
  \item \textbf{\texttt{post}}: registra no log se o \textit{pipeline} terminou com sucesso ou falha.
\end{itemize}

Com isso, o laboratório passou a ter um fluxo de \textbf{entrega contínua} completo até o ambiente de
testes. Para torná-lo \textit{deploy} contínuo em produção, bastaria acrescentar um estágio que
publicasse a imagem em um \textit{registry} e atualizasse os servidores de produção (por exemplo,
executando o Ansible); para mantê-lo como entrega contínua, esse estágio teria um \texttt{input} de
aprovação manual.

\subsubsection{Evidências da execução}

\printtela{jenkins-pipeline-sucesso}{Pipeline executado com sucesso e seus stages}%
  {Página do \textit{job} de \textit{pipeline} com a visualização dos estágios (\textit{Stage View}
   ou \textit{Pipeline Graph View}) mostrando todos os \textit{stages} em verde.}

\printtela{jenkins-pipeline-logs}{Logs do build do pipeline (Console Output)}%
  {\textit{Console Output} do \textit{build} do \textit{pipeline} mostrando a execução dos
   \textit{stages} (por exemplo, a saída do Maven ou do \texttt{sonar-scanner}) e
   \texttt{Finished: SUCCESS}.}

\printtela{jenkins-jenkinsfile}{Jenkinsfile configurado no job de pipeline}%
  {Tela de configuração do \textit{job} (\textit{Configure} $\rightarrow$ \textit{Pipeline}) com o
   \texttt{Jenkinsfile} visível no campo \textit{Script} (ou o arquivo no repositório GitHub).}

\subsubsection{Como o Jenkins automatiza etapas manuais}

Sem o Jenkins, cada nova versão exigiria que alguém executasse à mão a mesma sequência de passos:
atualizar o código, compilar, rodar os testes, executar a análise, gerar a imagem, parar o container
antigo, subir o novo e testar a API. Além de demorado, esse processo é sujeito a erros -- esquecer um
passo, executar na ordem errada, usar a versão errada do Java -- e depende da disponibilidade de uma
pessoa. Com o \textit{pipeline}, essas etapas rodam sempre da mesma forma, na mesma ordem, a cada
alteração; o processo para no primeiro erro, registra tudo em log, guarda o histórico de cada
\textit{build} e dá \textit{feedback} rápido para quem fez a mudança. É isso que permite entregar
com frequência e com confiança.

% =====================================================================
%  2.7 Integração entre as ferramentas
% =====================================================================
\subsection{Integração entre as ferramentas}\label{sec:integracao}

Cada ferramenta foi estudada separadamente, mas o objetivo final era entender como elas se combinam
em uma esteira DevOps (Figura~\ref{fig:integracao} e Tabela~\ref{tab:fases-ferramentas}).

\begin{figure}[H]
  \centering
  \caption{Integração entre as ferramentas no laboratório}\label{fig:integracao}
  \begin{tikzpicture}
    \tikzset{bloco/.style={caixa, text width=2.4cm, minimum height=1.2cm}}
    % Linha de cima: fluxo de CI/CD
    \node[caixacinza, text width=2.5cm, minimum height=1.2cm] (dev) at (0,0) {\textbf{Desenvolvedor}\\{\scriptsize edita e faz \textit{commit}}};
    \node[caixalaranja, text width=2.4cm, minimum height=1.2cm] (git) at (4.1,0) {\textbf{Git / GitHub}\\{\scriptsize versiona código e IaC}};
    \node[caixaescura, text width=2.4cm, minimum height=1.2cm] (jen) at (8.4,0) {Jenkins\\{\normalfont\scriptsize orquestra o \textit{pipeline}}};
    \node[bloco] (son) at (12.8,1.05) {\textbf{SonarQube}\\{\scriptsize qualidade e \textit{quality gate}}};
    \node[bloco] (doc) at (12.8,-1.05) {\textbf{Docker}\\{\scriptsize imagem e container}};
    % Linha de baixo: infraestrutura configurada
    \node[bloco] (ans) at (4.1,-3.1) {\textbf{Ansible}\\{\scriptsize \textit{control-node}}};
    \node[bloco] (app) at (8.4,-3.1) {\textbf{app01}\\{\scriptsize Java + Easy Notes}};
    \node[bloco] (db)  at (12.8,-3.1) {\textbf{db01}\\{\scriptsize MariaDB}};
    % Setas
    \draw[seta] (dev) -- node[rotulo, above=1pt]{\texttt{git push}} (git);
    \draw[seta] (git) -- node[rotulo, above=1pt]{\textit{poll}} node[rotulo, below=1pt]{checkout} (jen);
    \draw[seta] (jen.east) -- node[rotulo, above, sloped]{análise} (son.west);
    \draw[seta] (jen.east) -- node[rotulo, below, sloped]{\textit{build}/\textit{deploy}} (doc.west);
    \draw[seta, dashed] (git.south) -- node[rotulo, right=2pt]{\textit{playbooks}} (ans.north);
    \draw[seta] (ans) -- node[rotulo, above=1pt]{SSH} (app);
    \draw[seta] (ans.south) to[bend right=16] node[rotulo, above=1pt]{SSH} (db.south);
    \draw[seta, color=corverde] (app) -- node[rotulo, above=1pt]{JDBC 3306} (db);
    \draw[seta, dashed, color=corverde] (jen.south) to[out=-90, in=150]
         node[rotulo, pos=0.45, below left=0pt]{testes (JDBC)} (db.north west);
    \draw[seta, dashed, color=corverde] (doc.south) -- node[rotulo, right=2pt]{JDBC 3306} (db.north);
    % Base: Vagrant
    \coordinate (bl) at ($(dev.west |- ans.south)+(0,-1.65)$);
    \coordinate (br) at ($(db.east |- ans.south)+(0,-1.05)$);
    \draw[draw=corprimaria, fill=corprimaria!10, rounded corners=3pt] (bl) rectangle (br);
    \node[font=\scriptsize\bfseries, text=corprimaria] at ($(bl)!0.5!(br)$)
         {Vagrant + VMware: todas as VMs definidas em \texttt{Vagrantfile} (Infraestrutura como Código)};
  \end{tikzpicture}
  \fonte{Elaborado pelo autor (2026).}
\end{figure}

\begin{table}[H]
  \centering
  \caption{Etapas do ciclo DevOps e ferramentas utilizadas no laboratório}\label{tab:fases-ferramentas}
  \small
  \renewcommand{\arraystretch}{1.18}
  \begin{tabularx}{\linewidth}{>{\raggedright\arraybackslash}p{3.3cm} l >{\raggedright\arraybackslash}X}
    \toprule
    \textbf{Etapa} & \textbf{Ferramenta} & \textbf{Como foi usada} \\
    \midrule
    Codificar e versionar & Git e GitHub & Versionamento do código e de toda a configuração do laboratório. \\
    Provisionar a infraestrutura & Vagrant + VMware & Criação reprodutível das VMs a partir de arquivos de texto. \\
    Configurar servidores & Ansible & Preparação do sistema, MariaDB no \texttt{db01} e \textit{deploy} da aplicação no \texttt{app01}. \\
    \textit{Build} e testes & Maven (no Jenkins) & Compilação, empacotamento e testes a cada alteração. \\
    Qualidade e segurança & SonarQube & Análise estática e Quality Gate antes do empacotamento. \\
    Empacotar e \textit{deploy} & Docker + Jenkins & Imagem da aplicação e troca automática do container. \\
    Operar e monitorar & \texttt{systemd}, logs, \texttt{curl} & Verificação dos serviços e testes da API. \\
    \bottomrule
  \end{tabularx}
  \fonte{Elaborado pelo autor (2026).}
\end{table}

Na prática, o caminho de uma alteração fica assim: o desenvolvedor faz \cod|git push|; o Jenkins
detecta o novo \textit{commit}, compila e testa a aplicação contra o banco do \texttt{db01} --
configurado pelo Ansible em uma VM criada pelo Vagrant --; o SonarQube analisa o código e o Quality
Gate decide se ele pode seguir; se aprovado, o Docker empacota a aplicação e o Jenkins a coloca no ar.
Se algo falha, o \textit{pipeline} para e o \textit{feedback} chega rápido. O ponto central é que
\textbf{tudo é código} versionado no Git -- as máquinas, a configuração, a imagem e o próprio processo
de entrega --, o que torna o ambiente reprodutível e fácil de evoluir.

% =====================================================================
%  3 Conclusão
% =====================================================================
\section{Conclusão}\label{sec:conclusao}

As atividades práticas da disciplina de Fundamentos de DevOps permitiram construir, de forma
incremental, um laboratório que reproduz em pequena escala o ambiente de uma empresa. O Vagrant,
junto com o VMware, tornou possível criar e recriar sete máquinas virtuais a partir de arquivos de
texto; o Git e o GitHub garantiram o versionamento de todo o trabalho; o Ansible automatizou a
configuração dos servidores de banco de dados e de aplicação e o \textit{deploy} da aplicação Easy
Notes; o Docker empacotou essa aplicação em uma imagem portátil; o SonarQube avaliou a qualidade e a
segurança do código; e o Jenkins encadeou essas etapas em um \textit{pipeline} de CI/CD.

Mais do que aprender comandos, as práticas mostraram os princípios por trás do DevOps:
\textbf{automação} no lugar de tarefas manuais repetitivas, \textbf{Infraestrutura como Código} e
\textbf{versionamento} para ter ambientes reprodutíveis e rastreáveis, \textbf{idempotência} para
executar a automação com segurança quantas vezes for preciso, \textbf{qualidade contínua} com a
análise estática e o Quality Gate, e \textbf{\textit{feedback} rápido} por meio do \textit{pipeline}.
Os problemas enfrentados no caminho -- autenticação SSH entre as VMs, indentação do YAML,
permissões de diretórios, versões do Java, portas e \textit{firewall} -- também fizeram parte do
aprendizado, pois exigiram diagnosticar o ambiente camada por camada, como acontece no dia a dia de
operações.

Outro aprendizado importante foi o cuidado com informações sensíveis: senhas, chaves privadas e
\textit{tokens} não devem aparecer no código versionado nem nos prints, e ferramentas como o Ansible
Vault e as \textit{credentials} do Jenkins existem justamente para isso. Como evolução do laboratório,
seria interessante proteger as variáveis com o Ansible Vault, fazer o Jenkins acionar os
\textit{playbooks} do Ansible, publicar as imagens em um \textit{registry}, adicionar monitoramento
(por exemplo, com Prometheus e Grafana) e levar a mesma estrutura para um ambiente de nuvem. Conclui-se
que as ferramentas estudadas, quando integradas, transformam a entrega de software em um processo
padronizado, rápido e confiável, que é exatamente o que a cultura DevOps propõe.

% =====================================================================
%  4 Pergunta reflexiva final
% =====================================================================
% ATENÇÃO: esta resposta precisa estar com as SUAS palavras (exigência do professor).
% Use o texto abaixo como rascunho: revise, troque exemplos e acrescente sua experiência real.
\section{Pergunta reflexiva final}\label{sec:reflexao}

\begin{tcolorbox}[enhanced, colback=corprimaria!5, colframe=corprimaria, boxrule=0pt, leftrule=3pt,
    arc=0pt, outer arc=0pt, left=8pt, right=8pt, top=6pt, bottom=6pt]
  \itshape Qual das tecnologias estudadas na disciplina você ainda não conhecia ou conhecia pouco?
  Como a disciplina ajudou você a entender melhor o funcionamento do DevOps?
\end{tcolorbox}

Das tecnologias estudadas, a que eu menos conhecia era o \textbf{Ansible} e, junto com ele, o próprio
conceito de Infraestrutura como Código. Eu já usava Git e GitHub nos projetos da faculdade, mas
sempre tinha configurado servidores e ambientes ``na mão'', seguindo tutoriais e repetindo comandos
sem muito controle do que tinha sido feito. O Vagrant também era novidade para mim. Com o Ansible,
aprendi que é possível descrever em arquivos YAML como um servidor deve ficar e deixar a ferramenta
garantir esse estado, e que, graças à idempotência, executar o mesmo \textit{playbook} de novo não
estraga nada: ele só confirma que está tudo certo ou corrige o que estiver diferente.

A maior dificuldade foi fazer a comunicação entre as máquinas funcionar. O primeiro
\cod|ansible -m ping all| falhou com \textit{Permission denied}, e eu precisei entender de verdade
como funciona a autenticação por chave SSH: a chave privada fica no \textit{control-node} e a chave
pública vai para o \texttt{authorized\_keys} dos servidores. Também tive problemas com a indentação
do YAML e com o arquivo de configuração do Ansible, que era ignorado por estar em uma pasta com
permissão de escrita para todos. Resolver esses erros seguindo uma ordem de diagnóstico (máquina,
porta, \textit{firewall}, serviço, logs) me ensinou mais do que se tudo tivesse funcionado de
primeira.

A prática foi o que fez a teoria ter sentido. Ver o \texttt{db01} e o \texttt{app01} sendo
configurados sozinhos, a API do Easy Notes respondendo e a nota aparecendo gravada no banco mostrou
na prática o que significa automação. E, quando as ferramentas foram juntadas, ficou claro como elas
se conectam na cultura DevOps: o Vagrant cria a infraestrutura, o Git versiona tudo, o Ansible
configura os servidores, o Docker empacota a aplicação, o SonarQube verifica a qualidade e o Jenkins
orquestra todas essas etapas a cada \textit{commit}. Percebi que nenhuma delas, sozinha, ``é o
DevOps'': o DevOps está na integração entre elas e na forma de trabalhar, com automação, colaboração e
\textit{feedback} rápido.

Para a minha formação, esse conhecimento é muito importante. Hoje, mesmo quem trabalha principalmente
com desenvolvimento precisa entender como o código chega até a produção, e o mercado procura
profissionais que saibam automatizar, versionar e entregar software com qualidade e segurança. Como
já venho estudando Azure por conta própria, consegui relacionar o que vi na disciplina com a nuvem,
onde os mesmos conceitos de Infraestrutura como Código, \textit{pipelines} e containers aparecem o
tempo todo. A disciplina me deu uma base prática que pretendo continuar desenvolvendo nos próximos
projetos.


% ---------------------------------------------------------------------
%  Elementos pós-textuais
% ---------------------------------------------------------------------
% =====================================================================
%  Referências (formato ABNT, ordem alfabética)
% =====================================================================
\clearpage
\phantomsection
\addcontentsline{toc}{section}{Referências}
\renewcommand{\refname}{Referências}
\begingroup
\setstretch{1.0}
\small
\begin{thebibliography}{99}
  \setlength{\itemsep}{0.55em}

  \bibitem[CHACON; STRAUB(2014)]{chacon2014}
  CHACON, Scott; STRAUB, Ben. \textbf{Pro Git}. 2. ed. Berkeley: Apress, 2014. Disponível em:
  \url{https://git-scm.com/book/pt-br/v2}. Acesso em: 17 nov. 2026.

  \bibitem[DOCKER(2026)]{docker2026}
  DOCKER INC. \textbf{Docker Docs}: Dockerfile reference e Docker Engine. 2026. Disponível em:
  \url{https://docs.docker.com/}. Acesso em: 17 nov. 2026.

  \bibitem[HASHICORP(2026)]{hashicorp2026}
  HASHICORP. \textbf{Vagrant Documentation}. 2026. Disponível em:
  \url{https://developer.hashicorp.com/vagrant/docs}. Acesso em: 17 nov. 2026.

  \bibitem[HUMBLE; FARLEY(2014)]{humble2014}
  HUMBLE, Jez; FARLEY, David. \textbf{Entrega contínua}: como entregar software de forma rápida e
  confiável. Porto Alegre: Bookman, 2014.

  \bibitem[JENKINS(2026)]{jenkins2026}
  JENKINS PROJECT. \textbf{Jenkins User Documentation}: Installing Jenkins on Linux e Pipeline
  Syntax. 2026. Disponível em: \url{https://www.jenkins.io/doc/}. Acesso em: 17 nov. 2026.

  \bibitem[KIM et al.(2018)]{kim2018}
  KIM, Gene; HUMBLE, Jez; DEBOIS, Patrick; WILLIS, John. \textbf{Manual de DevOps}: como obter
  agilidade, confiabilidade e segurança em organizações tecnológicas. Rio de Janeiro: Alta Books,
  2018.

  \bibitem[MUNIZ(2026)]{muniz2026}
  MUNIZ, Augusto da Rosa. \textbf{Fundamentos de DevOps}: material das aulas 1 a 10 (Vagrant,
  Git, Ansible, Docker e SonarQube). Pelotas: Centro Universitário Senac RS -- UNISENAC, 2026.
  Material de aula (slides).

  \bibitem[RED HAT(2026)]{redhat2026}
  RED HAT. \textbf{Ansible Documentation}: playbooks, roles, inventory e variable precedence. 2026.
  Disponível em: \url{https://docs.ansible.com/}. Acesso em: 17 nov. 2026.

  \bibitem[SINGH(2021)]{singh2021}
  SINGH, Rajeev Kumar. \textbf{spring-boot-mysql-rest-api-tutorial} (Easy Notes). GitHub, 2021.
  Disponível em: \url{https://github.com/callicoder/spring-boot-mysql-rest-api-tutorial}.
  Acesso em: 17 nov. 2026.

  \bibitem[SONARSOURCE(2026)]{sonar2026}
  SONARSOURCE. \textbf{SonarQube Server Documentation}: SonarScanner CLI, quality gates e
  quality profiles. 2026. Disponível em: \url{https://docs.sonarsource.com/sonarqube-server/}.
  Acesso em: 17 nov. 2026.

\end{thebibliography}
\endgroup


\end{document}
